% ================================================================
% O singură sursă, două PDF-uri (compilează fiecare de două ori):
%
%   La clasă (interactiv, soluția se dezvăluie pe bucăți):
%     pdflatex -jobname=lesson-interactive lesson.tex
%
%   Compact (toți pașii într-un singur slide, pentru elevi / README):
%     pdflatex -jobname=lesson "\PassOptionsToClass{handout}{beamer}% ================================================================
% O singură sursă, două PDF-uri (compilează fiecare de două ori):
%
%   La clasă (interactiv, soluția se dezvăluie pe bucăți):
%     pdflatex -jobname=lesson-interactive lesson.tex
%
%   Compact (toți pașii într-un singur slide, pentru elevi / README):
%     pdflatex -jobname=lesson "\PassOptionsToClass{handout}{beamer}% ================================================================
% O singură sursă, două PDF-uri (compilează fiecare de două ori):
%
%   La clasă (interactiv, soluția se dezvăluie pe bucăți):
%     pdflatex -jobname=lesson-interactive lesson.tex
%
%   Compact (toți pașii într-un singur slide, pentru elevi / README):
%     pdflatex -jobname=lesson "\PassOptionsToClass{handout}{beamer}% ================================================================
% O singură sursă, două PDF-uri (compilează fiecare de două ori):
%
%   La clasă (interactiv, soluția se dezvăluie pe bucăți):
%     pdflatex -jobname=lesson-interactive lesson.tex
%
%   Compact (toți pașii într-un singur slide, pentru elevi / README):
%     pdflatex -jobname=lesson "\PassOptionsToClass{handout}{beamer}\input{lesson.tex}"
% ================================================================
\documentclass[9pt]{beamer}
\usepackage{../../../resources/themes/algolymp-dense}

% Evidențiere a doi termeni care se urmăresc prin substituții
\newcommand{\hla}[1]{\textbf{\textcolor{advAlert}{#1}}}
\newcommand{\hlb}[1]{\textbf{\textcolor{advAccent}{#1}}}

% Mod interactiv: întrebarea apare prima, apoi soluția se dezvăluie pe bucăți
% (\item<2->, \uncover<3->{...}, \begin{alertblock}<4->{...}); hint-ul se
% folosește doar unde elevii chiar trebuie împinși să se gândească.
% Totul se afișează cumulativ: ce apare la pasul 2 rămâne vizibil după.
\newenvironment<>{intrebare}{%
  \begin{uncoverenv}#1%
  \begin{block}{Întrebare}}{\end{block}\end{uncoverenv}}
\newenvironment<>{hint}{%
  \begin{uncoverenv}#1%
  \setbeamercolor{block title}{bg=advAlert!15,fg=advAlert}%
  \setbeamercolor{block body}{bg=advAlert!5,fg=black}%
  \begin{block}{Hint}}{\end{block}\end{uncoverenv}}

% Celulă de tabel vizibilă de la pasul #1
\newcommand{\at}[2]{\uncover<#1->{#2}}

\title{Algoritmul lui Euclid \textbar\ Exponențiere rapidă}
\subtitle{}
\author{Andrei Mihai Ciobanu}
\institute{Algolymp Summer School}
\date{}
\titlegraphic{\includegraphics[width=3cm]{../../../resources/algolymp-logo.png}}
\setbeamertemplate{frame footer}{Algolymp Summer School}

% ================================================================
\begin{document}

\begin{frame}\titlepage\end{frame}

\begin{frame}{Cuprins}\tableofcontents\end{frame}

% ================================================================
\section{Algoritmul lui Euclid}

\begin{frame}{Cel mai mare divizor comun}
  \begin{block}{Cel mai mare divizor comun}
    $\text{cmmdc}(a, b)$ = cel mai mare număr natural $d$ astfel încât $d \mid a$ și $d \mid b$.
    Se mai notează și $\gcd(a, b)$ (\textit{greatest common divisor}) sau doar $(a, b)$.
  \end{block}
  \begin{block}{Cel mai mic multiplu comun}
    $\text{cmmmc}(a, b)$ = cel mai mic număr natural $m$ astfel încât $a \mid m$ și $b \mid m$.
    Se mai notează și $\text{lcm}(a, b)$ (\textit{lowest common multiple}) sau doar $[a, b]$.
  \end{block}
  \smallskip
  Exemple:
  \begin{itemize}
    \item $\gcd(12, 8) = 4$ \quad $\gcd(17, 5) = 1$ \quad $\gcd(100, 75) = 25$
    \item $\text{lcm}(4, 6) = 12$ \quad $\text{lcm}(3, 5) = 15$
  \end{itemize}
\end{frame}

\begin{frame}{Calculul cmmdc}
  \begin{intrebare}
    Cum se calculează $\gcd(84, 36)$? Dar $\gcd(10^{18},\ 6 \cdot 10^{17})$?
  \end{intrebare}
  \begin{itemize}
    \item<2-> Se listează divizorii lui $a$ în $O(\sqrt{a})$ și se testează fiecare pe $b$.
      Pentru $\gcd(84, 36) = 12$ metoda funcționează.
    \item<2-> Pentru numere de ordinul $10^{18}$ rezultă circa $10^9$ operații \textemdash{} prea lent.
  \end{itemize}
  \begin{alertblock}<3->{Necesitatea unei metode mai rapide}
    \textbf{Algoritmul lui Euclid}: sub $100$ de pași pentru numere de $10^{18}$.
  \end{alertblock}
\end{frame}

\begin{frame}{Ideea de bază: $\gcd(a, b) = \gcd(a - b,\ b)$}
  \begin{intrebare}
    Fie $d$ un divizor comun al lui $a$ și $b$. Ce se poate afirma despre $a - b$?
  \end{intrebare}
  \begin{hint}<2->
    Se scriu $a = d \cdot x$ și $b = d \cdot y$, apoi se calculează $a - b$.
  \end{hint}
  \begin{itemize}
    \item<3-> $a - b = d \cdot (x - y)$, deci $d \mid a - b$.
    \item<3-> Reciproc, dacă $d \mid a - b$ și $d \mid b$, atunci $d \mid (a - b) + b = a$.
  \end{itemize}
  \begin{alertblock}<4->{Concluzie}
    Perechile $(a, b)$ și $(a - b,\ b)$ au \textbf{aceiași divizori comuni}, deci
    $\gcd(a, b) = \gcd(a - b,\ b)$.
  \end{alertblock}
\end{frame}

\begin{frame}{Scăderi repetate}
  Regula $\gcd(a, b) = \gcd(a - b, b)$ se aplică repetat până la cazul de bază $\gcd(a, 0) = a$.
  \begin{intrebare}
    Cât este $\gcd(12, 8)$ dacă se folosesc \textbf{doar} scăderi?
  \end{intrebare}
  \uncover<2->{
    \[
      \gcd(12, 8) = \gcd(4, 8) = \gcd(4, 4) = \gcd(0, 4) = 4 \quad \checkmark
    \]
  }
  \begin{intrebare}<3->
    Câți pași sunt necesari pentru $\gcd(1000, 1)$?
  \end{intrebare}
  \uncover<4->{
    \[
      \gcd(1000, 1) \rightarrow \gcd(999, 1) \rightarrow \gcd(998, 1) \rightarrow \cdots
      \rightarrow \gcd(0, 1) = 1
    \]
    \textbf{1000 de pași} pentru un rezultat evident \textemdash{} procedeul este prea lent pentru valori mari.
  }
\end{frame}

\begin{frame}{De la scăderi la rest}
  \begin{intrebare}
    Cum pot fi înlocuite scăderile repetate ale lui $b$ cu o singură operație?
  \end{intrebare}
  \uncover<2->{
    Prin împărțirea cu rest, $a = q \cdot b + r$, cu $0 \leq r < b$. Restul $r$ se notează
    $a \bmod b$ (restul împărțirii lui $a$ la $b$):
    \[
      \gcd(a, b) = \gcd(a - b, b) = \gcd(a - 2b, b) = \cdots = \gcd(a \bmod b,\ b)
    \]
  }
  \begin{alertblock}<3->{Algoritmul lui Euclid}
    $\gcd(a, b) = \gcd(b,\ a \bmod b)$, \quad cu cazul de bază $\gcd(a, 0) = a$.
  \end{alertblock}
  \uncover<4->{Exemplu: $\gcd(48, 18)$}
  \begin{itemize}
    \item<4-> $\gcd(48, 18) = \gcd(18, 12)$ \quad ($48 \bmod 18 = 12$)
    \item<4-> $\gcd(18, 12) = \gcd(12, 6)$ \quad\ ($18 \bmod 12 = 6$)
    \item<5-> $\gcd(12,\ 6) = \gcd(6,\ 0)$ \quad\ \ ($12 \bmod 6 = 0$)
    \item<5-> $\gcd(6, 0) = 6$ \checkmark
  \end{itemize}
\end{frame}

\begin{frame}{Proprietăți}
  \begin{itemize}
    \item $\gcd(a, 0) = \gcd(a, a) = |a|$ \quad (\textit{cazuri de bază})
    \item $\gcd(a, b) = \gcd(b, a)$ \quad (\textit{comutativitate})
    \item $\gcd(a, \gcd(b, c)) = \gcd(\gcd(a, b), c)$ \quad (\textit{asociativitate})
    \item $\gcd(k \cdot a, k \cdot b) = k \cdot \gcd(a, b)$, pentru $k > 0$
    \item dacă $\gcd(a, b) = d$, atunci $\gcd\!\left(\dfrac{a}{d},\ \dfrac{b}{d}\right) = 1$
    \item $a \cdot b = \gcd(a, b) \cdot \text{lcm}(a, b)$
    \item dacă $\gcd(a, b) = 1$, se spune că $a$ și $b$ sunt \textit{prime între ele} (\textit{coprime})
    \item $\gcd(a, b) = \gcd(a \pm k \cdot b,\ b)$, $\forall\, k \in \mathbb{N}$ \quad
      $\Rightarrow$ \quad $\gcd(a, b) = \gcd(b,\ a \bmod b)$
    \item \textbf{Prin descompunere în factori primi:} dacă
      $n = \prod p_i^{a_i}$ și $m = \prod p_i^{b_i}$, atunci
      $\gcd(n, m) = \prod p_i^{\min(a_i,\, b_i)}$ și
      $\text{lcm}(n, m) = \prod p_i^{\max(a_i,\, b_i)}$ \\
      Exemplu: $n = 2^3 \cdot 3^2,\ m = 2^1 \cdot 3^4 \;\Rightarrow\;
        \gcd = 2^1 \cdot 3^2 = 18,\;
        \text{lcm} = 2^3 \cdot 3^4 = 648$
  \end{itemize}
\end{frame}

\begin{frame}[fragile]{Algoritmul lui Euclid \textemdash{} implementare iterativă}
  \begin{block}{Explicația implementării}
    \begin{itemize}
      \item Se pornește de la $\gcd(a, b) = \gcd(a \bmod b,\ b) = \gcd(b,\ a \bmod b)$
      \item Cât timp $b > 0$, se calculează restul $r = a \bmod b$;
        apoi noul $a$ devine $b$, iar noul $b$ devine $r$
      \item După fiecare pas, noul $b$ este mai mic decât noul $a$ ($r < b$), deci primul
        termen rămâne cel mai mare și operația $a \bmod b$ se poate efectua în continuare
      \item Dacă inițial $a < b$, atunci $a \bmod b = a$, iar primul pas doar interschimbă
        valorile: $(a, b) \rightarrow (b, a)$
      \item Când $b = 0$, rezultatul este $a$ (cazul de bază $\gcd(a, 0) = a$)
    \end{itemize}
  \end{block}
  \begin{uncoverenv}<2->
    \begin{lstlisting}
int gcd(int a, int b) {
    while (b != 0) {
        int r = a % b;
        a = b;
        b = r;
    }
    return a;
}
    \end{lstlisting}
    \begin{itemize}
      \item Timp: $O(\log \min(a, b))$ \quad Spațiu: $O(1)$
    \end{itemize}
  \end{uncoverenv}
\end{frame}

\begin{frame}[fragile]{Cel mai mic multiplu comun \textemdash{} calcul}
  \begin{intrebare}
    Cum se calculează $\text{lcm}(a, b)$ dacă se cunoaște $\gcd(a, b)$?
  \end{intrebare}
  \uncover<2->{
    Din $a \cdot b = \gcd(a, b) \cdot \text{lcm}(a, b)$ rezultă
    \lstinline[keepspaces=true]|lcm = (a / gcd(a, b)) * b|.
  }
  \begin{alertblock}<3->{Ordinea operațiilor}
    Este indicat ca împărțirea să se facă prima: $a \cdot b$ poate produce overflow chiar dacă
    $\text{lcm}(a, b)$ încape în tipul de date, iar astfel rezultatul rămâne corect.
  \end{alertblock}
\end{frame}

\begin{frame}{Complexitatea algoritmului lui Euclid}
  \begin{intrebare}
    Cât devine $a$ după \textbf{doi} pași ai algoritmului?
  \end{intrebare}
  \begin{hint}<2->
    Se separă cazurile $b \leq a/2$ și $b > a/2$. Cât este $a \bmod b$ în fiecare caz?
  \end{hint}
  \begin{itemize}
    \item<3-> dacă $b \leq a/2$: $a \bmod b < b \leq a/2$
    \item<3-> dacă $b > a/2$: $a \bmod b = a - b < a/2$
  \end{itemize}
  \begin{alertblock}<4->{Concluzie}
    În ambele cazuri $a \bmod b < a/2$, deci după \textbf{doi} pași cel mai mare număr
    se înjumătățește cel puțin $\Rightarrow$ $O(\log \min(a, b))$ pași.
  \end{alertblock}
  \beginnote{$\log$ = numărul de înjumătățiri succesive până când un număr ajunge la $1$; $\log_2 10^{18} \approx 60$}
\end{frame}

% ================================================================
\section{Algoritmul lui Euclid extins}

\begin{frame}{Identitatea lui Bézout}
  \begin{block}{Identitatea lui Bézout}
    Pentru orice $a, b \in \mathbb{Z}$, există $x, y \in \mathbb{Z}$ astfel încât:
    \[ a \cdot x + b \cdot y = \gcd(a, b) \]
  \end{block}
  Exemplu: $22 \cdot 3 + 13 \cdot (-5) = 66 - 65 = 1 = \gcd(22, 13)$.
  \begin{intrebare}
    Cum se determină și numerele $x, y$, nu doar $\gcd$?
  \end{intrebare}
  \uncover<2->{
    Algoritmul lui Euclid reduce $(a, b)$ la o pereche mai mică, $(b,\ a \bmod b)$.
    Soluția pentru perechea mai mică poate fi folosită pentru perechea inițială?
  }
  \beginnote{Secțiunea conține un concept ușor mai dificil \textemdash{} poate fi reluată după parcurgerea restului lecției.}
\end{frame}

\begin{frame}{Euclid extins \textemdash{} ideea}
  Fie $g = \gcd(a, b)$ și ecuația $a \cdot x + b \cdot y = g$.
  Se scrie $a = q \cdot b + r$, unde $q = \lfloor a / b \rfloor$ și $r = a \bmod b$.
  \begin{intrebare}
    Ce formă are ecuația după înlocuirea lui $a$ și gruparea termenilor care conțin $b$?
  \end{intrebare}
  \uncover<2->{
    \[
      (\underbrace{q \cdot b + r}_{a}) \cdot x + b \cdot y = g
      \;\;\Rightarrow\;\;
      b \cdot (q \cdot x + y) + r \cdot x = g
    \]
  }
  \uncover<3->{
    Cu notațiile $x' = q \cdot x + y$ și $y' = x$, ecuația devine $b \cdot x' + r \cdot y' = g$.
  }
  \begin{itemize}
    \item<4-> Forma este aceeași, cu $(b,\ r)$ în loc de $(a,\ b)$ \textemdash{} exact un pas din algoritmul lui Euclid.
    \item<4-> Reducerea continuă până la cazul de bază $b = 0$: $g \cdot x + 0 \cdot y = g \;\Rightarrow\; x = 1,\ y = 0$
  \end{itemize}
  \begin{alertblock}<5->{Reconstrucție la întoarcere}
    Din $x' = q \cdot x + y$ și $y' = x$ rezultă $x = y'$, \quad $y = x' - q \cdot y'$.
  \end{alertblock}
\end{frame}

\begin{frame}{Exemplu pe hârtie: $22x + 13y = \gcd(22, 13)$}
  \textbf{Pasul 1: Se aplică algoritmul lui Euclid}
  \begin{itemize}
    \item $22 = 1 \cdot 13 + 9$ \hfill {\footnotesize\textcolor{gray}{$\Rightarrow 9 = 22 - 1 \cdot 13$}}
    \item $13 = 1 \cdot 9 + 4$ \hfill {\footnotesize\textcolor{gray}{$\Rightarrow 4 = 13 - 1 \cdot 9$}}
    \item $9 = 2 \cdot 4 + 1$ \hfill {\footnotesize\textcolor{gray}{$\Rightarrow 1 = 9 - 2 \cdot 4$}}
    \item $4 = 4 \cdot 1 + 0 \;\Rightarrow\; \gcd(22, 13) = 1$
  \end{itemize}
  \begin{intrebare}
    Cum se exprimă $1$ folosind doar $22$ și $13$?
  \end{intrebare}
  \uncover<2->{\textbf{Pasul 2: Substituție în ordine inversă, de jos în sus}}
  \begin{itemize}
    \item<2-> $1 = 1 \cdot \hla{9} - 2 \cdot \hlb{4}$
    \item<3-> $1 = 1 \cdot \hla{9} - 2 \cdot (\hlb{13} - 1 \cdot \hla{9}) = 3 \cdot \hla{9} - 2 \cdot \hlb{13}$
    \item<4-> $1 = 3 \cdot (1 \cdot \hla{22} - 1 \cdot \hlb{13}) - 2 \cdot \hlb{13} = 3 \cdot \hla{22} - 5 \cdot \hlb{13}$
  \end{itemize}
  \begin{alertblock}<5->{Soluție}
    $x = 3,\ y = -5$ \quad \checkmark \quad ($22 \cdot 3 + 13 \cdot (-5) = 66 - 65 = 1$)
  \end{alertblock}
\end{frame}

\begin{frame}[fragile]{Euclid extins \textemdash{} implementare iterativă}
  \begin{lstlisting}
// Gaseste x, y cu: a*x + b*y = gcd(a, b)
int gcd_ext(int a, int b, int &x, int &y) {
    int x0 = 1, y0 = 0;
    int x1 = 0, y1 = 1;
    while (b != 0) {
        int q  = a / b;
        int xn = x0 - q * x1;
        int yn = y0 - q * y1;
        int r  = a % b;
        a = b;  b = r;
        x0 = x1; y0 = y1;
        x1 = xn; y1 = yn;
    }
    x = x0; y = y0;
    return a;
}
  \end{lstlisting}
  \begin{itemize}
    \item \lstinline|int &x| = parametru transmis prin referință: modificările din funcție sunt vizibile în apelant
    \item Timp: $O(\log \min(a, b))$ \quad Spațiu: $O(1)$
  \end{itemize}
\end{frame}

\begin{frame}{Inițializarea coeficienților $x_0, y_0, x_1, y_1$}
  Fie $a_0, b_0$ valorile inițiale. La orice pas al buclei sunt menținute relațiile:
  \[
    a = a_0 \cdot x_0 + b_0 \cdot y_0 \qquad b = a_0 \cdot x_1 + b_0 \cdot y_1
  \]
  \begin{intrebare}
    Ce valori trebuie să aibă $x_0, y_0, x_1, y_1$ înainte de primul pas?
  \end{intrebare}
  \begin{itemize}
    \item<2-> Înainte de primul pas, $a = a_0$ și $b = b_0$.
    \item<2-> $a_0 = a_0 \cdot 1 + b_0 \cdot 0$ și $b_0 = a_0 \cdot 0 + b_0 \cdot 1$,
      deci $x_0 = 1,\ y_0 = 0,\ x_1 = 0,\ y_1 = 1$.
    \item<3-> La final $b = 0$ și $a = \gcd$, deci $x_0, y_0$ sunt coeficienții căutați.
  \end{itemize}
\end{frame}

% ================================================================
\section{Exponențiere rapidă}

\begin{frame}{Exponențiere rapidă \textemdash{} ideea}
  \begin{block}{Problema}
    Calculul eficient al lui $a^n$. Varianta naivă ($n$ înmulțiri) are complexitate $O(n)$.
  \end{block}
  \begin{intrebare}
    Cum se poate calcula $2^{10}$ cu mai puțin de $9$ înmulțiri?
  \end{intrebare}
  \begin{itemize}
    \item<2-> $2^{10} = (2^5)^2$: $2^5$ se calculează o singură dată și se ridică la pătrat
    \item<2-> Aceeași jumătate nu se mai calculează de două ori: exponentul se înjumătățește
  \end{itemize}
  \begin{intrebare}<3->
    Cum se procedează dacă exponentul este impar, de exemplu $2^5$?
  \end{intrebare}
  \begin{itemize}
    \item<4-> Se scoate o bază în afară: $2^5 = 2 \cdot 2^4$, iar $2^4$ are exponent par
  \end{itemize}
  \uncover<4->{
    \[
      a^n = \begin{cases}
        1 & n = 0 \\
        \left(a^{n/2}\right)^2 & n \text{ par} \\
        a \cdot a^{n-1} & n \text{ impar}
      \end{cases}
    \]
  }
\end{frame}

\begin{frame}{Numărul de pași}
  Exemplu: $2^{10} \rightarrow 2^5 \rightarrow 2^4 \rightarrow 2^2 \rightarrow 2^1 \rightarrow 2^0$
  \quad (5 pași în loc de 10)
  \begin{intrebare}
    Câți pași sunt necesari pentru $n = 10^9$?
  \end{intrebare}
  \begin{itemize}
    \item<2-> $10^9$ se împarte la $2$ de circa $30$ de ori, deoarece $2^{30} \approx 10^9$
    \item<2-> La fiecare înjumătățire se face cel mult o înmulțire suplimentară pentru exponentul impar,
      deci cel mult circa $2 \cdot 30$ înmulțiri
  \end{itemize}
  \begin{alertblock}<3->{Concluzie}
    Complexitate $O(\log n)$, în loc de $10^9$ pași.
  \end{alertblock}
\end{frame}

\begin{frame}[fragile]{Exponențiere rapidă \textemdash{} implementare iterativă}
  \begin{lstlisting}
long long power(long long base, int exp) {
    long long result = 1;
    while (exp > 0) {
        if (exp % 2 == 1)
            result *= base;
        base *= base;
        exp /= 2;
    }
    return result;
}
  \end{lstlisting}
  \begin{itemize}
    \item Timp: $O(\log n)$ \quad Spațiu: $O(1)$
    \item \alert{Atenție:} \lstinline|base * base| poate da overflow pentru baze mari \textemdash{}
      în concursuri se lucrează de obicei modulo un număr prim
  \end{itemize}
\end{frame}

\begin{frame}{Urmărirea execuției: \texttt{power(3, 5)}}
  \begin{intrebare}
    Ce valori au \texttt{result}, \texttt{base} și \texttt{exp} la finalul fiecărei iterații?
  \end{intrebare}
  \begin{center}
    \begin{tabular}{cccc}
      \toprule
      Iterația & \texttt{result} & \texttt{base} & \texttt{exp} \\
      \midrule
      început & $1$ & $3$ & $5$ \\
      \at{2}{$1$ (\texttt{exp} $= 5$, impar)} & \at{2}{$3$} & \at{2}{$9$} & \at{2}{$2$} \\
      \at{3}{$2$ (\texttt{exp} $= 2$, par)}   & \at{3}{$3$} & \at{3}{$81$} & \at{3}{$1$} \\
      \at{4}{$3$ (\texttt{exp} $= 1$, impar)} & \at{4}{$243$} & \at{4}{$6561$} & \at{4}{$0$} \\
      \bottomrule
    \end{tabular}
  \end{center}
  \begin{alertblock}<5->{Rezultat}
    \texttt{power(3, 5)} $= 243 = 3^5$ \quad \checkmark
  \end{alertblock}
\end{frame}

\begin{frame}{Legătura cu scrierea binară}
  În exemplul \texttt{power(3, 5)}: $5 = 101_{(2)}$, iar \texttt{base} a luat valorile $3^1,\ 3^2,\ 3^4$.
  \begin{intrebare}
    Ce legătură există între algoritm și scrierea lui $n$ în baza $2$?
  \end{intrebare}
  \begin{hint}<2->
    Se descompune $5 = 4 + 1$ și se scrie $3^5$ ca produs de puteri ale lui $3$.
  \end{hint}
  \begin{itemize}
    \item<3-> $3^5 = 3^4 \cdot 3^1$; \texttt{base} parcurge $a^{2^0}, a^{2^1}, a^{2^2}, \ldots$, iar
      \texttt{result} se înmulțește doar cu bazele corespunzătoare \textbf{biților egali cu $1$}
      din scrierea binară a lui $n$
  \end{itemize}
  \begin{alertblock}<4->{Concluzie}
    Numărul de biți ai lui $n$ este $\approx \log_2 n$, deci $O(\log n)$ iterații.
  \end{alertblock}
\end{frame}

% ================================================================
\section{Bonus \textemdash{} Implementări recursive}

\begin{frame}{Recursivitatea}
  \begin{block}{Definiție}
    O funcție este \textbf{recursivă} dacă se apelează pe ea însăși.
  \end{block}
  Orice funcție recursivă are nevoie de:
  \begin{itemize}
    \item Un \textbf{caz de bază} \textemdash{} condiția de oprire
    \item Un \textbf{apel recursiv} pe o subproblemă mai mică
  \end{itemize}
  Exemplu: $n! = n \cdot (n-1)!$, cu cazul de bază $0! = 1$
  \begin{intrebare}
    Ce se întâmplă dacă lipsește cazul de bază?
  \end{intrebare}
  \begin{itemize}
    \item<2-> Funcția s-ar apela la nesfârșit până se umple stiva de apeluri \textemdash{}
      \textbf{stack overflow}
  \end{itemize}
  \begin{block}<3->{\textcolor{advAlert}{Atenție} \textemdash{} dezavantajele recursivității}
    Overhead de apeluri de funcție (ușor mai lent) și risc de \textbf{stack overflow} pentru intrări
    foarte mari. Codul recursiv reflectă direct definiția matematică, dar în concursuri
    este preferată varianta iterativă.
  \end{block}
\end{frame}

\begin{frame}[fragile]{Euclid \textemdash{} implementare recursivă}
  \begin{intrebare}
    Cum se scrie $\gcd$ recursiv pornind de la $\gcd(a, b) = \gcd(b,\ a \bmod b)$?
    Care este cazul de bază?
  \end{intrebare}
  \begin{uncoverenv}<2->
    \begin{lstlisting}
int gcd(int a, int b) {
    if (b == 0)
        return a;
    return gcd(b, a % b);
}
    \end{lstlisting}
    \begin{itemize}
      \item Cazul de bază este $\gcd(a, 0) = a$
      \item Codul reflectă direct proprietatea $\gcd(a, b) = \gcd(b,\ a \bmod b)$
      \item Aceeași complexitate ca varianta iterativă: $O(\log \min(a, b))$
    \end{itemize}
  \end{uncoverenv}
\end{frame}

\begin{frame}[fragile]{Euclid extins \textemdash{} implementare recursivă}
  \begin{intrebare}
    Ce valori au $x$ și $y$ în cazul de bază $b = 0$? Cum se reconstruiesc la întoarcere?
  \end{intrebare}
  \begin{itemize}
    \item<2-> Cazul de bază: $a \cdot x + 0 \cdot y = a$, deci $x = 1,\ y = 0$
    \item<2-> La întoarcere: $x = y'$, \quad $y = x' - q \cdot y'$
  \end{itemize}
  \begin{uncoverenv}<3->
    \begin{lstlisting}
int gcd_ext(int a, int b, int &x, int &y) {
    if (b == 0) {
        x = 1; y = 0;
        return a;
    }
    int x1, y1;
    int d = gcd_ext(b, a % b, x1, y1);
    x = y1;
    y = x1 - (a / b) * y1;
    return d;
}
    \end{lstlisting}
    \begin{itemize}
      \item Timp: $O(\log \min(a, b))$ \quad Spațiu: $O(\log \min(a, b))$ \textemdash{} stiva de apeluri
    \end{itemize}
  \end{uncoverenv}
\end{frame}

\begin{frame}[fragile]{Exponențiere rapidă \textemdash{} implementare recursivă}
  \begin{lstlisting}
long long power(long long base, int exp) {
    if (exp == 0)
        return 1;
    if (exp % 2 == 0) {
        long long half = power(base, exp / 2);
        return half * half;
    }
    return base * power(base, exp - 1);
}
  \end{lstlisting}
  \begin{itemize}
    \item Timp: $O(\log n)$ \quad Spațiu: $O(\log n)$ \textemdash{} stiva de apeluri
  \end{itemize}
  \begin{intrebare}
    De ce nu se scrie direct
    \lstinline[keepspaces=true]|return power(base, exp / 2) * power(base, exp / 2);|?
  \end{intrebare}
  \begin{itemize}
    \item<2-> Două apeluri la fiecare nivel $\Rightarrow$ numărul de apeluri se dublează la fiecare nivel
    \item<2-> În total $O(n)$ apeluri \textemdash{} câștigul se pierde; cu \lstinline|half * half|
      se face \textbf{un singur} apel pe nivel: $O(\log n)$
  \end{itemize}
\end{frame}

% ================================================================
\section{Recapitulare}

\begin{frame}{Recapitulare \textemdash{} Complexități}
  \begin{center}
  \begin{tabular}{lll}
    \toprule
    Algoritm & Timp & Spațiu \\
    \midrule
    $\gcd$ iterativ                & $O(\log \min(a,b))$ & $O(1)$ \\
    $\gcd$ extins iterativ         & $O(\log \min(a,b))$ & $O(1)$ \\
    $\gcd$ recursiv                & $O(\log \min(a,b))$ & $O(\log \min(a,b))$ \\
    $\gcd$ extins recursiv         & $O(\log \min(a,b))$ & $O(\log \min(a,b))$ \\
    Exponențiere rapidă iterativă  & $O(\log n)$         & $O(1)$ \\
    Exponențiere rapidă recursivă  & $O(\log n)$         & $O(\log n)$ \\
    \bottomrule
  \end{tabular}
  \end{center}
\end{frame}

\begin{frame}{Probleme de antrenament}
  \begin{block}{AlgoCode}
    \begin{enumerate}
      \item \href{https://code.algolymp.com/problems/1557}{\textbf{GCD and LCM}}
            \textcolor{gray}{(AlgoCode \#1557)}
      \item \href{https://code.algolymp.com/problems/4230}{\textbf{GCD and LCM Queries}}
            \textcolor{gray}{(AlgoCode \#4230)}
      \item \href{https://code.algolymp.com/problems/4270}{\textbf{GCD Power Queries}}
            \textcolor{gray}{(AlgoCode \#4270)}
    \end{enumerate}
  \end{block}
  \begin{block}{Altele}
    \begin{enumerate}
      \item \href{https://www.pbinfo.ro/probleme/58/cmmdc}{\textbf{cmmdc}}
            \textcolor{gray}{(pbinfo.ro)}
      \item \href{https://www.pbinfo.ro/probleme/59/cmmmc}{\textbf{cmmmc}}
            \textcolor{gray}{(pbinfo.ro)}
      \item \href{https://www.infoarena.ro/problema/euclid3}{\textbf{euclid3}}
            \textcolor{gray}{(infoarena.ro)}
      \item \href{https://infoarena.ro/problema/lgput}{\textbf{lgput}}
            \textcolor{gray}{(infoarena.ro)}
    \end{enumerate}
  \end{block}
\end{frame}

\begin{frame}{Resurse}
  \begin{enumerate}
    \item \href{https://www.infoarena.ro/algoritmul-lui-euclid}{\textbf{infoarena.ro}}
          \textemdash{} Algoritmul lui Euclid (RO)
    \item \href{https://cp-algorithms.com/algebra/euclid-algorithm.html}{\textbf{cp-algorithms.com}}
          \textemdash{} Euclidean algorithm (EN)
    \item \href{https://cp-algorithms.com/algebra/extended-euclid-algorithm.html}{\textbf{cp-algorithms.com}}
          \textemdash{} Extended Euclidean algorithm (EN)
    \item \href{https://cp-algorithms.com/algebra/binary-exp.html}{\textbf{cp-algorithms.com}}
          \textemdash{} Binary exponentiation (EN)
    \item \href{https://www.pbinfo.ro/articole/3873/recursivitate}{\textbf{pbinfo.ro}}
          \textemdash{} Recursivitate (RO)
    \item \href{https://cppreference.com}{\textbf{cppreference.com}}
          \textemdash{} documentație completă C++
  \end{enumerate}
\end{frame}

\end{document}
"
% ================================================================
\documentclass[9pt]{beamer}
\usepackage{../../../resources/themes/algolymp-dense}

% Evidențiere a doi termeni care se urmăresc prin substituții
\newcommand{\hla}[1]{\textbf{\textcolor{advAlert}{#1}}}
\newcommand{\hlb}[1]{\textbf{\textcolor{advAccent}{#1}}}

% Mod interactiv: întrebarea apare prima, apoi soluția se dezvăluie pe bucăți
% (\item<2->, \uncover<3->{...}, \begin{alertblock}<4->{...}); hint-ul se
% folosește doar unde elevii chiar trebuie împinși să se gândească.
% Totul se afișează cumulativ: ce apare la pasul 2 rămâne vizibil după.
\newenvironment<>{intrebare}{%
  \begin{uncoverenv}#1%
  \begin{block}{Întrebare}}{\end{block}\end{uncoverenv}}
\newenvironment<>{hint}{%
  \begin{uncoverenv}#1%
  \setbeamercolor{block title}{bg=advAlert!15,fg=advAlert}%
  \setbeamercolor{block body}{bg=advAlert!5,fg=black}%
  \begin{block}{Hint}}{\end{block}\end{uncoverenv}}

% Celulă de tabel vizibilă de la pasul #1
\newcommand{\at}[2]{\uncover<#1->{#2}}

\title{Algoritmul lui Euclid \textbar\ Exponențiere rapidă}
\subtitle{}
\author{Andrei Mihai Ciobanu}
\institute{Algolymp Summer School}
\date{}
\titlegraphic{\includegraphics[width=3cm]{../../../resources/algolymp-logo.png}}
\setbeamertemplate{frame footer}{Algolymp Summer School}

% ================================================================
\begin{document}

\begin{frame}\titlepage\end{frame}

\begin{frame}{Cuprins}\tableofcontents\end{frame}

% ================================================================
\section{Algoritmul lui Euclid}

\begin{frame}{Cel mai mare divizor comun}
  \begin{block}{Cel mai mare divizor comun}
    $\text{cmmdc}(a, b)$ = cel mai mare număr natural $d$ astfel încât $d \mid a$ și $d \mid b$.
    Se mai notează și $\gcd(a, b)$ (\textit{greatest common divisor}) sau doar $(a, b)$.
  \end{block}
  \begin{block}{Cel mai mic multiplu comun}
    $\text{cmmmc}(a, b)$ = cel mai mic număr natural $m$ astfel încât $a \mid m$ și $b \mid m$.
    Se mai notează și $\text{lcm}(a, b)$ (\textit{lowest common multiple}) sau doar $[a, b]$.
  \end{block}
  \smallskip
  Exemple:
  \begin{itemize}
    \item $\gcd(12, 8) = 4$ \quad $\gcd(17, 5) = 1$ \quad $\gcd(100, 75) = 25$
    \item $\text{lcm}(4, 6) = 12$ \quad $\text{lcm}(3, 5) = 15$
  \end{itemize}
\end{frame}

\begin{frame}{Calculul cmmdc}
  \begin{intrebare}
    Cum se calculează $\gcd(84, 36)$? Dar $\gcd(10^{18},\ 6 \cdot 10^{17})$?
  \end{intrebare}
  \begin{itemize}
    \item<2-> Se listează divizorii lui $a$ în $O(\sqrt{a})$ și se testează fiecare pe $b$.
      Pentru $\gcd(84, 36) = 12$ metoda funcționează.
    \item<2-> Pentru numere de ordinul $10^{18}$ rezultă circa $10^9$ operații \textemdash{} prea lent.
  \end{itemize}
  \begin{alertblock}<3->{Necesitatea unei metode mai rapide}
    \textbf{Algoritmul lui Euclid}: sub $100$ de pași pentru numere de $10^{18}$.
  \end{alertblock}
\end{frame}

\begin{frame}{Ideea de bază: $\gcd(a, b) = \gcd(a - b,\ b)$}
  \begin{intrebare}
    Fie $d$ un divizor comun al lui $a$ și $b$. Ce se poate afirma despre $a - b$?
  \end{intrebare}
  \begin{hint}<2->
    Se scriu $a = d \cdot x$ și $b = d \cdot y$, apoi se calculează $a - b$.
  \end{hint}
  \begin{itemize}
    \item<3-> $a - b = d \cdot (x - y)$, deci $d \mid a - b$.
    \item<3-> Reciproc, dacă $d \mid a - b$ și $d \mid b$, atunci $d \mid (a - b) + b = a$.
  \end{itemize}
  \begin{alertblock}<4->{Concluzie}
    Perechile $(a, b)$ și $(a - b,\ b)$ au \textbf{aceiași divizori comuni}, deci
    $\gcd(a, b) = \gcd(a - b,\ b)$.
  \end{alertblock}
\end{frame}

\begin{frame}{Scăderi repetate}
  Regula $\gcd(a, b) = \gcd(a - b, b)$ se aplică repetat până la cazul de bază $\gcd(a, 0) = a$.
  \begin{intrebare}
    Cât este $\gcd(12, 8)$ dacă se folosesc \textbf{doar} scăderi?
  \end{intrebare}
  \uncover<2->{
    \[
      \gcd(12, 8) = \gcd(4, 8) = \gcd(4, 4) = \gcd(0, 4) = 4 \quad \checkmark
    \]
  }
  \begin{intrebare}<3->
    Câți pași sunt necesari pentru $\gcd(1000, 1)$?
  \end{intrebare}
  \uncover<4->{
    \[
      \gcd(1000, 1) \rightarrow \gcd(999, 1) \rightarrow \gcd(998, 1) \rightarrow \cdots
      \rightarrow \gcd(0, 1) = 1
    \]
    \textbf{1000 de pași} pentru un rezultat evident \textemdash{} procedeul este prea lent pentru valori mari.
  }
\end{frame}

\begin{frame}{De la scăderi la rest}
  \begin{intrebare}
    Cum pot fi înlocuite scăderile repetate ale lui $b$ cu o singură operație?
  \end{intrebare}
  \uncover<2->{
    Prin împărțirea cu rest, $a = q \cdot b + r$, cu $0 \leq r < b$. Restul $r$ se notează
    $a \bmod b$ (restul împărțirii lui $a$ la $b$):
    \[
      \gcd(a, b) = \gcd(a - b, b) = \gcd(a - 2b, b) = \cdots = \gcd(a \bmod b,\ b)
    \]
  }
  \begin{alertblock}<3->{Algoritmul lui Euclid}
    $\gcd(a, b) = \gcd(b,\ a \bmod b)$, \quad cu cazul de bază $\gcd(a, 0) = a$.
  \end{alertblock}
  \uncover<4->{Exemplu: $\gcd(48, 18)$}
  \begin{itemize}
    \item<4-> $\gcd(48, 18) = \gcd(18, 12)$ \quad ($48 \bmod 18 = 12$)
    \item<4-> $\gcd(18, 12) = \gcd(12, 6)$ \quad\ ($18 \bmod 12 = 6$)
    \item<5-> $\gcd(12,\ 6) = \gcd(6,\ 0)$ \quad\ \ ($12 \bmod 6 = 0$)
    \item<5-> $\gcd(6, 0) = 6$ \checkmark
  \end{itemize}
\end{frame}

\begin{frame}{Proprietăți}
  \begin{itemize}
    \item $\gcd(a, 0) = \gcd(a, a) = |a|$ \quad (\textit{cazuri de bază})
    \item $\gcd(a, b) = \gcd(b, a)$ \quad (\textit{comutativitate})
    \item $\gcd(a, \gcd(b, c)) = \gcd(\gcd(a, b), c)$ \quad (\textit{asociativitate})
    \item $\gcd(k \cdot a, k \cdot b) = k \cdot \gcd(a, b)$, pentru $k > 0$
    \item dacă $\gcd(a, b) = d$, atunci $\gcd\!\left(\dfrac{a}{d},\ \dfrac{b}{d}\right) = 1$
    \item $a \cdot b = \gcd(a, b) \cdot \text{lcm}(a, b)$
    \item dacă $\gcd(a, b) = 1$, se spune că $a$ și $b$ sunt \textit{prime între ele} (\textit{coprime})
    \item $\gcd(a, b) = \gcd(a \pm k \cdot b,\ b)$, $\forall\, k \in \mathbb{N}$ \quad
      $\Rightarrow$ \quad $\gcd(a, b) = \gcd(b,\ a \bmod b)$
    \item \textbf{Prin descompunere în factori primi:} dacă
      $n = \prod p_i^{a_i}$ și $m = \prod p_i^{b_i}$, atunci
      $\gcd(n, m) = \prod p_i^{\min(a_i,\, b_i)}$ și
      $\text{lcm}(n, m) = \prod p_i^{\max(a_i,\, b_i)}$ \\
      Exemplu: $n = 2^3 \cdot 3^2,\ m = 2^1 \cdot 3^4 \;\Rightarrow\;
        \gcd = 2^1 \cdot 3^2 = 18,\;
        \text{lcm} = 2^3 \cdot 3^4 = 648$
  \end{itemize}
\end{frame}

\begin{frame}[fragile]{Algoritmul lui Euclid \textemdash{} implementare iterativă}
  \begin{block}{Explicația implementării}
    \begin{itemize}
      \item Se pornește de la $\gcd(a, b) = \gcd(a \bmod b,\ b) = \gcd(b,\ a \bmod b)$
      \item Cât timp $b > 0$, se calculează restul $r = a \bmod b$;
        apoi noul $a$ devine $b$, iar noul $b$ devine $r$
      \item După fiecare pas, noul $b$ este mai mic decât noul $a$ ($r < b$), deci primul
        termen rămâne cel mai mare și operația $a \bmod b$ se poate efectua în continuare
      \item Dacă inițial $a < b$, atunci $a \bmod b = a$, iar primul pas doar interschimbă
        valorile: $(a, b) \rightarrow (b, a)$
      \item Când $b = 0$, rezultatul este $a$ (cazul de bază $\gcd(a, 0) = a$)
    \end{itemize}
  \end{block}
  \begin{uncoverenv}<2->
    \begin{lstlisting}
int gcd(int a, int b) {
    while (b != 0) {
        int r = a % b;
        a = b;
        b = r;
    }
    return a;
}
    \end{lstlisting}
    \begin{itemize}
      \item Timp: $O(\log \min(a, b))$ \quad Spațiu: $O(1)$
    \end{itemize}
  \end{uncoverenv}
\end{frame}

\begin{frame}[fragile]{Cel mai mic multiplu comun \textemdash{} calcul}
  \begin{intrebare}
    Cum se calculează $\text{lcm}(a, b)$ dacă se cunoaște $\gcd(a, b)$?
  \end{intrebare}
  \uncover<2->{
    Din $a \cdot b = \gcd(a, b) \cdot \text{lcm}(a, b)$ rezultă
    \lstinline[keepspaces=true]|lcm = (a / gcd(a, b)) * b|.
  }
  \begin{alertblock}<3->{Ordinea operațiilor}
    Este indicat ca împărțirea să se facă prima: $a \cdot b$ poate produce overflow chiar dacă
    $\text{lcm}(a, b)$ încape în tipul de date, iar astfel rezultatul rămâne corect.
  \end{alertblock}
\end{frame}

\begin{frame}{Complexitatea algoritmului lui Euclid}
  \begin{intrebare}
    Cât devine $a$ după \textbf{doi} pași ai algoritmului?
  \end{intrebare}
  \begin{hint}<2->
    Se separă cazurile $b \leq a/2$ și $b > a/2$. Cât este $a \bmod b$ în fiecare caz?
  \end{hint}
  \begin{itemize}
    \item<3-> dacă $b \leq a/2$: $a \bmod b < b \leq a/2$
    \item<3-> dacă $b > a/2$: $a \bmod b = a - b < a/2$
  \end{itemize}
  \begin{alertblock}<4->{Concluzie}
    În ambele cazuri $a \bmod b < a/2$, deci după \textbf{doi} pași cel mai mare număr
    se înjumătățește cel puțin $\Rightarrow$ $O(\log \min(a, b))$ pași.
  \end{alertblock}
  \beginnote{$\log$ = numărul de înjumătățiri succesive până când un număr ajunge la $1$; $\log_2 10^{18} \approx 60$}
\end{frame}

% ================================================================
\section{Algoritmul lui Euclid extins}

\begin{frame}{Identitatea lui Bézout}
  \begin{block}{Identitatea lui Bézout}
    Pentru orice $a, b \in \mathbb{Z}$, există $x, y \in \mathbb{Z}$ astfel încât:
    \[ a \cdot x + b \cdot y = \gcd(a, b) \]
  \end{block}
  Exemplu: $22 \cdot 3 + 13 \cdot (-5) = 66 - 65 = 1 = \gcd(22, 13)$.
  \begin{intrebare}
    Cum se determină și numerele $x, y$, nu doar $\gcd$?
  \end{intrebare}
  \uncover<2->{
    Algoritmul lui Euclid reduce $(a, b)$ la o pereche mai mică, $(b,\ a \bmod b)$.
    Soluția pentru perechea mai mică poate fi folosită pentru perechea inițială?
  }
  \beginnote{Secțiunea conține un concept ușor mai dificil \textemdash{} poate fi reluată după parcurgerea restului lecției.}
\end{frame}

\begin{frame}{Euclid extins \textemdash{} ideea}
  Fie $g = \gcd(a, b)$ și ecuația $a \cdot x + b \cdot y = g$.
  Se scrie $a = q \cdot b + r$, unde $q = \lfloor a / b \rfloor$ și $r = a \bmod b$.
  \begin{intrebare}
    Ce formă are ecuația după înlocuirea lui $a$ și gruparea termenilor care conțin $b$?
  \end{intrebare}
  \uncover<2->{
    \[
      (\underbrace{q \cdot b + r}_{a}) \cdot x + b \cdot y = g
      \;\;\Rightarrow\;\;
      b \cdot (q \cdot x + y) + r \cdot x = g
    \]
  }
  \uncover<3->{
    Cu notațiile $x' = q \cdot x + y$ și $y' = x$, ecuația devine $b \cdot x' + r \cdot y' = g$.
  }
  \begin{itemize}
    \item<4-> Forma este aceeași, cu $(b,\ r)$ în loc de $(a,\ b)$ \textemdash{} exact un pas din algoritmul lui Euclid.
    \item<4-> Reducerea continuă până la cazul de bază $b = 0$: $g \cdot x + 0 \cdot y = g \;\Rightarrow\; x = 1,\ y = 0$
  \end{itemize}
  \begin{alertblock}<5->{Reconstrucție la întoarcere}
    Din $x' = q \cdot x + y$ și $y' = x$ rezultă $x = y'$, \quad $y = x' - q \cdot y'$.
  \end{alertblock}
\end{frame}

\begin{frame}{Exemplu pe hârtie: $22x + 13y = \gcd(22, 13)$}
  \textbf{Pasul 1: Se aplică algoritmul lui Euclid}
  \begin{itemize}
    \item $22 = 1 \cdot 13 + 9$ \hfill {\footnotesize\textcolor{gray}{$\Rightarrow 9 = 22 - 1 \cdot 13$}}
    \item $13 = 1 \cdot 9 + 4$ \hfill {\footnotesize\textcolor{gray}{$\Rightarrow 4 = 13 - 1 \cdot 9$}}
    \item $9 = 2 \cdot 4 + 1$ \hfill {\footnotesize\textcolor{gray}{$\Rightarrow 1 = 9 - 2 \cdot 4$}}
    \item $4 = 4 \cdot 1 + 0 \;\Rightarrow\; \gcd(22, 13) = 1$
  \end{itemize}
  \begin{intrebare}
    Cum se exprimă $1$ folosind doar $22$ și $13$?
  \end{intrebare}
  \uncover<2->{\textbf{Pasul 2: Substituție în ordine inversă, de jos în sus}}
  \begin{itemize}
    \item<2-> $1 = 1 \cdot \hla{9} - 2 \cdot \hlb{4}$
    \item<3-> $1 = 1 \cdot \hla{9} - 2 \cdot (\hlb{13} - 1 \cdot \hla{9}) = 3 \cdot \hla{9} - 2 \cdot \hlb{13}$
    \item<4-> $1 = 3 \cdot (1 \cdot \hla{22} - 1 \cdot \hlb{13}) - 2 \cdot \hlb{13} = 3 \cdot \hla{22} - 5 \cdot \hlb{13}$
  \end{itemize}
  \begin{alertblock}<5->{Soluție}
    $x = 3,\ y = -5$ \quad \checkmark \quad ($22 \cdot 3 + 13 \cdot (-5) = 66 - 65 = 1$)
  \end{alertblock}
\end{frame}

\begin{frame}[fragile]{Euclid extins \textemdash{} implementare iterativă}
  \begin{lstlisting}
// Gaseste x, y cu: a*x + b*y = gcd(a, b)
int gcd_ext(int a, int b, int &x, int &y) {
    int x0 = 1, y0 = 0;
    int x1 = 0, y1 = 1;
    while (b != 0) {
        int q  = a / b;
        int xn = x0 - q * x1;
        int yn = y0 - q * y1;
        int r  = a % b;
        a = b;  b = r;
        x0 = x1; y0 = y1;
        x1 = xn; y1 = yn;
    }
    x = x0; y = y0;
    return a;
}
  \end{lstlisting}
  \begin{itemize}
    \item \lstinline|int &x| = parametru transmis prin referință: modificările din funcție sunt vizibile în apelant
    \item Timp: $O(\log \min(a, b))$ \quad Spațiu: $O(1)$
  \end{itemize}
\end{frame}

\begin{frame}{Inițializarea coeficienților $x_0, y_0, x_1, y_1$}
  Fie $a_0, b_0$ valorile inițiale. La orice pas al buclei sunt menținute relațiile:
  \[
    a = a_0 \cdot x_0 + b_0 \cdot y_0 \qquad b = a_0 \cdot x_1 + b_0 \cdot y_1
  \]
  \begin{intrebare}
    Ce valori trebuie să aibă $x_0, y_0, x_1, y_1$ înainte de primul pas?
  \end{intrebare}
  \begin{itemize}
    \item<2-> Înainte de primul pas, $a = a_0$ și $b = b_0$.
    \item<2-> $a_0 = a_0 \cdot 1 + b_0 \cdot 0$ și $b_0 = a_0 \cdot 0 + b_0 \cdot 1$,
      deci $x_0 = 1,\ y_0 = 0,\ x_1 = 0,\ y_1 = 1$.
    \item<3-> La final $b = 0$ și $a = \gcd$, deci $x_0, y_0$ sunt coeficienții căutați.
  \end{itemize}
\end{frame}

% ================================================================
\section{Exponențiere rapidă}

\begin{frame}{Exponențiere rapidă \textemdash{} ideea}
  \begin{block}{Problema}
    Calculul eficient al lui $a^n$. Varianta naivă ($n$ înmulțiri) are complexitate $O(n)$.
  \end{block}
  \begin{intrebare}
    Cum se poate calcula $2^{10}$ cu mai puțin de $9$ înmulțiri?
  \end{intrebare}
  \begin{itemize}
    \item<2-> $2^{10} = (2^5)^2$: $2^5$ se calculează o singură dată și se ridică la pătrat
    \item<2-> Aceeași jumătate nu se mai calculează de două ori: exponentul se înjumătățește
  \end{itemize}
  \begin{intrebare}<3->
    Cum se procedează dacă exponentul este impar, de exemplu $2^5$?
  \end{intrebare}
  \begin{itemize}
    \item<4-> Se scoate o bază în afară: $2^5 = 2 \cdot 2^4$, iar $2^4$ are exponent par
  \end{itemize}
  \uncover<4->{
    \[
      a^n = \begin{cases}
        1 & n = 0 \\
        \left(a^{n/2}\right)^2 & n \text{ par} \\
        a \cdot a^{n-1} & n \text{ impar}
      \end{cases}
    \]
  }
\end{frame}

\begin{frame}{Numărul de pași}
  Exemplu: $2^{10} \rightarrow 2^5 \rightarrow 2^4 \rightarrow 2^2 \rightarrow 2^1 \rightarrow 2^0$
  \quad (5 pași în loc de 10)
  \begin{intrebare}
    Câți pași sunt necesari pentru $n = 10^9$?
  \end{intrebare}
  \begin{itemize}
    \item<2-> $10^9$ se împarte la $2$ de circa $30$ de ori, deoarece $2^{30} \approx 10^9$
    \item<2-> La fiecare înjumătățire se face cel mult o înmulțire suplimentară pentru exponentul impar,
      deci cel mult circa $2 \cdot 30$ înmulțiri
  \end{itemize}
  \begin{alertblock}<3->{Concluzie}
    Complexitate $O(\log n)$, în loc de $10^9$ pași.
  \end{alertblock}
\end{frame}

\begin{frame}[fragile]{Exponențiere rapidă \textemdash{} implementare iterativă}
  \begin{lstlisting}
long long power(long long base, int exp) {
    long long result = 1;
    while (exp > 0) {
        if (exp % 2 == 1)
            result *= base;
        base *= base;
        exp /= 2;
    }
    return result;
}
  \end{lstlisting}
  \begin{itemize}
    \item Timp: $O(\log n)$ \quad Spațiu: $O(1)$
    \item \alert{Atenție:} \lstinline|base * base| poate da overflow pentru baze mari \textemdash{}
      în concursuri se lucrează de obicei modulo un număr prim
  \end{itemize}
\end{frame}

\begin{frame}{Urmărirea execuției: \texttt{power(3, 5)}}
  \begin{intrebare}
    Ce valori au \texttt{result}, \texttt{base} și \texttt{exp} la finalul fiecărei iterații?
  \end{intrebare}
  \begin{center}
    \begin{tabular}{cccc}
      \toprule
      Iterația & \texttt{result} & \texttt{base} & \texttt{exp} \\
      \midrule
      început & $1$ & $3$ & $5$ \\
      \at{2}{$1$ (\texttt{exp} $= 5$, impar)} & \at{2}{$3$} & \at{2}{$9$} & \at{2}{$2$} \\
      \at{3}{$2$ (\texttt{exp} $= 2$, par)}   & \at{3}{$3$} & \at{3}{$81$} & \at{3}{$1$} \\
      \at{4}{$3$ (\texttt{exp} $= 1$, impar)} & \at{4}{$243$} & \at{4}{$6561$} & \at{4}{$0$} \\
      \bottomrule
    \end{tabular}
  \end{center}
  \begin{alertblock}<5->{Rezultat}
    \texttt{power(3, 5)} $= 243 = 3^5$ \quad \checkmark
  \end{alertblock}
\end{frame}

\begin{frame}{Legătura cu scrierea binară}
  În exemplul \texttt{power(3, 5)}: $5 = 101_{(2)}$, iar \texttt{base} a luat valorile $3^1,\ 3^2,\ 3^4$.
  \begin{intrebare}
    Ce legătură există între algoritm și scrierea lui $n$ în baza $2$?
  \end{intrebare}
  \begin{hint}<2->
    Se descompune $5 = 4 + 1$ și se scrie $3^5$ ca produs de puteri ale lui $3$.
  \end{hint}
  \begin{itemize}
    \item<3-> $3^5 = 3^4 \cdot 3^1$; \texttt{base} parcurge $a^{2^0}, a^{2^1}, a^{2^2}, \ldots$, iar
      \texttt{result} se înmulțește doar cu bazele corespunzătoare \textbf{biților egali cu $1$}
      din scrierea binară a lui $n$
  \end{itemize}
  \begin{alertblock}<4->{Concluzie}
    Numărul de biți ai lui $n$ este $\approx \log_2 n$, deci $O(\log n)$ iterații.
  \end{alertblock}
\end{frame}

% ================================================================
\section{Bonus \textemdash{} Implementări recursive}

\begin{frame}{Recursivitatea}
  \begin{block}{Definiție}
    O funcție este \textbf{recursivă} dacă se apelează pe ea însăși.
  \end{block}
  Orice funcție recursivă are nevoie de:
  \begin{itemize}
    \item Un \textbf{caz de bază} \textemdash{} condiția de oprire
    \item Un \textbf{apel recursiv} pe o subproblemă mai mică
  \end{itemize}
  Exemplu: $n! = n \cdot (n-1)!$, cu cazul de bază $0! = 1$
  \begin{intrebare}
    Ce se întâmplă dacă lipsește cazul de bază?
  \end{intrebare}
  \begin{itemize}
    \item<2-> Funcția s-ar apela la nesfârșit până se umple stiva de apeluri \textemdash{}
      \textbf{stack overflow}
  \end{itemize}
  \begin{block}<3->{\textcolor{advAlert}{Atenție} \textemdash{} dezavantajele recursivității}
    Overhead de apeluri de funcție (ușor mai lent) și risc de \textbf{stack overflow} pentru intrări
    foarte mari. Codul recursiv reflectă direct definiția matematică, dar în concursuri
    este preferată varianta iterativă.
  \end{block}
\end{frame}

\begin{frame}[fragile]{Euclid \textemdash{} implementare recursivă}
  \begin{intrebare}
    Cum se scrie $\gcd$ recursiv pornind de la $\gcd(a, b) = \gcd(b,\ a \bmod b)$?
    Care este cazul de bază?
  \end{intrebare}
  \begin{uncoverenv}<2->
    \begin{lstlisting}
int gcd(int a, int b) {
    if (b == 0)
        return a;
    return gcd(b, a % b);
}
    \end{lstlisting}
    \begin{itemize}
      \item Cazul de bază este $\gcd(a, 0) = a$
      \item Codul reflectă direct proprietatea $\gcd(a, b) = \gcd(b,\ a \bmod b)$
      \item Aceeași complexitate ca varianta iterativă: $O(\log \min(a, b))$
    \end{itemize}
  \end{uncoverenv}
\end{frame}

\begin{frame}[fragile]{Euclid extins \textemdash{} implementare recursivă}
  \begin{intrebare}
    Ce valori au $x$ și $y$ în cazul de bază $b = 0$? Cum se reconstruiesc la întoarcere?
  \end{intrebare}
  \begin{itemize}
    \item<2-> Cazul de bază: $a \cdot x + 0 \cdot y = a$, deci $x = 1,\ y = 0$
    \item<2-> La întoarcere: $x = y'$, \quad $y = x' - q \cdot y'$
  \end{itemize}
  \begin{uncoverenv}<3->
    \begin{lstlisting}
int gcd_ext(int a, int b, int &x, int &y) {
    if (b == 0) {
        x = 1; y = 0;
        return a;
    }
    int x1, y1;
    int d = gcd_ext(b, a % b, x1, y1);
    x = y1;
    y = x1 - (a / b) * y1;
    return d;
}
    \end{lstlisting}
    \begin{itemize}
      \item Timp: $O(\log \min(a, b))$ \quad Spațiu: $O(\log \min(a, b))$ \textemdash{} stiva de apeluri
    \end{itemize}
  \end{uncoverenv}
\end{frame}

\begin{frame}[fragile]{Exponențiere rapidă \textemdash{} implementare recursivă}
  \begin{lstlisting}
long long power(long long base, int exp) {
    if (exp == 0)
        return 1;
    if (exp % 2 == 0) {
        long long half = power(base, exp / 2);
        return half * half;
    }
    return base * power(base, exp - 1);
}
  \end{lstlisting}
  \begin{itemize}
    \item Timp: $O(\log n)$ \quad Spațiu: $O(\log n)$ \textemdash{} stiva de apeluri
  \end{itemize}
  \begin{intrebare}
    De ce nu se scrie direct
    \lstinline[keepspaces=true]|return power(base, exp / 2) * power(base, exp / 2);|?
  \end{intrebare}
  \begin{itemize}
    \item<2-> Două apeluri la fiecare nivel $\Rightarrow$ numărul de apeluri se dublează la fiecare nivel
    \item<2-> În total $O(n)$ apeluri \textemdash{} câștigul se pierde; cu \lstinline|half * half|
      se face \textbf{un singur} apel pe nivel: $O(\log n)$
  \end{itemize}
\end{frame}

% ================================================================
\section{Recapitulare}

\begin{frame}{Recapitulare \textemdash{} Complexități}
  \begin{center}
  \begin{tabular}{lll}
    \toprule
    Algoritm & Timp & Spațiu \\
    \midrule
    $\gcd$ iterativ                & $O(\log \min(a,b))$ & $O(1)$ \\
    $\gcd$ extins iterativ         & $O(\log \min(a,b))$ & $O(1)$ \\
    $\gcd$ recursiv                & $O(\log \min(a,b))$ & $O(\log \min(a,b))$ \\
    $\gcd$ extins recursiv         & $O(\log \min(a,b))$ & $O(\log \min(a,b))$ \\
    Exponențiere rapidă iterativă  & $O(\log n)$         & $O(1)$ \\
    Exponențiere rapidă recursivă  & $O(\log n)$         & $O(\log n)$ \\
    \bottomrule
  \end{tabular}
  \end{center}
\end{frame}

\begin{frame}{Probleme de antrenament}
  \begin{block}{AlgoCode}
    \begin{enumerate}
      \item \href{https://code.algolymp.com/problems/1557}{\textbf{GCD and LCM}}
            \textcolor{gray}{(AlgoCode \#1557)}
      \item \href{https://code.algolymp.com/problems/4230}{\textbf{GCD and LCM Queries}}
            \textcolor{gray}{(AlgoCode \#4230)}
      \item \href{https://code.algolymp.com/problems/4270}{\textbf{GCD Power Queries}}
            \textcolor{gray}{(AlgoCode \#4270)}
    \end{enumerate}
  \end{block}
  \begin{block}{Altele}
    \begin{enumerate}
      \item \href{https://www.pbinfo.ro/probleme/58/cmmdc}{\textbf{cmmdc}}
            \textcolor{gray}{(pbinfo.ro)}
      \item \href{https://www.pbinfo.ro/probleme/59/cmmmc}{\textbf{cmmmc}}
            \textcolor{gray}{(pbinfo.ro)}
      \item \href{https://www.infoarena.ro/problema/euclid3}{\textbf{euclid3}}
            \textcolor{gray}{(infoarena.ro)}
      \item \href{https://infoarena.ro/problema/lgput}{\textbf{lgput}}
            \textcolor{gray}{(infoarena.ro)}
    \end{enumerate}
  \end{block}
\end{frame}

\begin{frame}{Resurse}
  \begin{enumerate}
    \item \href{https://www.infoarena.ro/algoritmul-lui-euclid}{\textbf{infoarena.ro}}
          \textemdash{} Algoritmul lui Euclid (RO)
    \item \href{https://cp-algorithms.com/algebra/euclid-algorithm.html}{\textbf{cp-algorithms.com}}
          \textemdash{} Euclidean algorithm (EN)
    \item \href{https://cp-algorithms.com/algebra/extended-euclid-algorithm.html}{\textbf{cp-algorithms.com}}
          \textemdash{} Extended Euclidean algorithm (EN)
    \item \href{https://cp-algorithms.com/algebra/binary-exp.html}{\textbf{cp-algorithms.com}}
          \textemdash{} Binary exponentiation (EN)
    \item \href{https://www.pbinfo.ro/articole/3873/recursivitate}{\textbf{pbinfo.ro}}
          \textemdash{} Recursivitate (RO)
    \item \href{https://cppreference.com}{\textbf{cppreference.com}}
          \textemdash{} documentație completă C++
  \end{enumerate}
\end{frame}

\end{document}
"
% ================================================================
\documentclass[9pt]{beamer}
\usepackage{../../../resources/themes/algolymp-dense}

% Evidențiere a doi termeni care se urmăresc prin substituții
\newcommand{\hla}[1]{\textbf{\textcolor{advAlert}{#1}}}
\newcommand{\hlb}[1]{\textbf{\textcolor{advAccent}{#1}}}

% Mod interactiv: întrebarea apare prima, apoi soluția se dezvăluie pe bucăți
% (\item<2->, \uncover<3->{...}, \begin{alertblock}<4->{...}); hint-ul se
% folosește doar unde elevii chiar trebuie împinși să se gândească.
% Totul se afișează cumulativ: ce apare la pasul 2 rămâne vizibil după.
\newenvironment<>{intrebare}{%
  \begin{uncoverenv}#1%
  \begin{block}{Întrebare}}{\end{block}\end{uncoverenv}}
\newenvironment<>{hint}{%
  \begin{uncoverenv}#1%
  \setbeamercolor{block title}{bg=advAlert!15,fg=advAlert}%
  \setbeamercolor{block body}{bg=advAlert!5,fg=black}%
  \begin{block}{Hint}}{\end{block}\end{uncoverenv}}

% Celulă de tabel vizibilă de la pasul #1
\newcommand{\at}[2]{\uncover<#1->{#2}}

\title{Algoritmul lui Euclid \textbar\ Exponențiere rapidă}
\subtitle{}
\author{Andrei Mihai Ciobanu}
\institute{Algolymp Summer School}
\date{}
\titlegraphic{\includegraphics[width=3cm]{../../../resources/algolymp-logo.png}}
\setbeamertemplate{frame footer}{Algolymp Summer School}

% ================================================================
\begin{document}

\begin{frame}\titlepage\end{frame}

\begin{frame}{Cuprins}\tableofcontents\end{frame}

% ================================================================
\section{Algoritmul lui Euclid}

\begin{frame}{Cel mai mare divizor comun}
  \begin{block}{Cel mai mare divizor comun}
    $\text{cmmdc}(a, b)$ = cel mai mare număr natural $d$ astfel încât $d \mid a$ și $d \mid b$.
    Se mai notează și $\gcd(a, b)$ (\textit{greatest common divisor}) sau doar $(a, b)$.
  \end{block}
  \begin{block}{Cel mai mic multiplu comun}
    $\text{cmmmc}(a, b)$ = cel mai mic număr natural $m$ astfel încât $a \mid m$ și $b \mid m$.
    Se mai notează și $\text{lcm}(a, b)$ (\textit{lowest common multiple}) sau doar $[a, b]$.
  \end{block}
  \smallskip
  Exemple:
  \begin{itemize}
    \item $\gcd(12, 8) = 4$ \quad $\gcd(17, 5) = 1$ \quad $\gcd(100, 75) = 25$
    \item $\text{lcm}(4, 6) = 12$ \quad $\text{lcm}(3, 5) = 15$
  \end{itemize}
\end{frame}

\begin{frame}{Calculul cmmdc}
  \begin{intrebare}
    Cum se calculează $\gcd(84, 36)$? Dar $\gcd(10^{18},\ 6 \cdot 10^{17})$?
  \end{intrebare}
  \begin{itemize}
    \item<2-> Se listează divizorii lui $a$ în $O(\sqrt{a})$ și se testează fiecare pe $b$.
      Pentru $\gcd(84, 36) = 12$ metoda funcționează.
    \item<2-> Pentru numere de ordinul $10^{18}$ rezultă circa $10^9$ operații \textemdash{} prea lent.
  \end{itemize}
  \begin{alertblock}<3->{Necesitatea unei metode mai rapide}
    \textbf{Algoritmul lui Euclid}: sub $100$ de pași pentru numere de $10^{18}$.
  \end{alertblock}
\end{frame}

\begin{frame}{Ideea de bază: $\gcd(a, b) = \gcd(a - b,\ b)$}
  \begin{intrebare}
    Fie $d$ un divizor comun al lui $a$ și $b$. Ce se poate afirma despre $a - b$?
  \end{intrebare}
  \begin{hint}<2->
    Se scriu $a = d \cdot x$ și $b = d \cdot y$, apoi se calculează $a - b$.
  \end{hint}
  \begin{itemize}
    \item<3-> $a - b = d \cdot (x - y)$, deci $d \mid a - b$.
    \item<3-> Reciproc, dacă $d \mid a - b$ și $d \mid b$, atunci $d \mid (a - b) + b = a$.
  \end{itemize}
  \begin{alertblock}<4->{Concluzie}
    Perechile $(a, b)$ și $(a - b,\ b)$ au \textbf{aceiași divizori comuni}, deci
    $\gcd(a, b) = \gcd(a - b,\ b)$.
  \end{alertblock}
\end{frame}

\begin{frame}{Scăderi repetate}
  Regula $\gcd(a, b) = \gcd(a - b, b)$ se aplică repetat până la cazul de bază $\gcd(a, 0) = a$.
  \begin{intrebare}
    Cât este $\gcd(12, 8)$ dacă se folosesc \textbf{doar} scăderi?
  \end{intrebare}
  \uncover<2->{
    \[
      \gcd(12, 8) = \gcd(4, 8) = \gcd(4, 4) = \gcd(0, 4) = 4 \quad \checkmark
    \]
  }
  \begin{intrebare}<3->
    Câți pași sunt necesari pentru $\gcd(1000, 1)$?
  \end{intrebare}
  \uncover<4->{
    \[
      \gcd(1000, 1) \rightarrow \gcd(999, 1) \rightarrow \gcd(998, 1) \rightarrow \cdots
      \rightarrow \gcd(0, 1) = 1
    \]
    \textbf{1000 de pași} pentru un rezultat evident \textemdash{} procedeul este prea lent pentru valori mari.
  }
\end{frame}

\begin{frame}{De la scăderi la rest}
  \begin{intrebare}
    Cum pot fi înlocuite scăderile repetate ale lui $b$ cu o singură operație?
  \end{intrebare}
  \uncover<2->{
    Prin împărțirea cu rest, $a = q \cdot b + r$, cu $0 \leq r < b$. Restul $r$ se notează
    $a \bmod b$ (restul împărțirii lui $a$ la $b$):
    \[
      \gcd(a, b) = \gcd(a - b, b) = \gcd(a - 2b, b) = \cdots = \gcd(a \bmod b,\ b)
    \]
  }
  \begin{alertblock}<3->{Algoritmul lui Euclid}
    $\gcd(a, b) = \gcd(b,\ a \bmod b)$, \quad cu cazul de bază $\gcd(a, 0) = a$.
  \end{alertblock}
  \uncover<4->{Exemplu: $\gcd(48, 18)$}
  \begin{itemize}
    \item<4-> $\gcd(48, 18) = \gcd(18, 12)$ \quad ($48 \bmod 18 = 12$)
    \item<4-> $\gcd(18, 12) = \gcd(12, 6)$ \quad\ ($18 \bmod 12 = 6$)
    \item<5-> $\gcd(12,\ 6) = \gcd(6,\ 0)$ \quad\ \ ($12 \bmod 6 = 0$)
    \item<5-> $\gcd(6, 0) = 6$ \checkmark
  \end{itemize}
\end{frame}

\begin{frame}{Proprietăți}
  \begin{itemize}
    \item $\gcd(a, 0) = \gcd(a, a) = |a|$ \quad (\textit{cazuri de bază})
    \item $\gcd(a, b) = \gcd(b, a)$ \quad (\textit{comutativitate})
    \item $\gcd(a, \gcd(b, c)) = \gcd(\gcd(a, b), c)$ \quad (\textit{asociativitate})
    \item $\gcd(k \cdot a, k \cdot b) = k \cdot \gcd(a, b)$, pentru $k > 0$
    \item dacă $\gcd(a, b) = d$, atunci $\gcd\!\left(\dfrac{a}{d},\ \dfrac{b}{d}\right) = 1$
    \item $a \cdot b = \gcd(a, b) \cdot \text{lcm}(a, b)$
    \item dacă $\gcd(a, b) = 1$, se spune că $a$ și $b$ sunt \textit{prime între ele} (\textit{coprime})
    \item $\gcd(a, b) = \gcd(a \pm k \cdot b,\ b)$, $\forall\, k \in \mathbb{N}$ \quad
      $\Rightarrow$ \quad $\gcd(a, b) = \gcd(b,\ a \bmod b)$
    \item \textbf{Prin descompunere în factori primi:} dacă
      $n = \prod p_i^{a_i}$ și $m = \prod p_i^{b_i}$, atunci
      $\gcd(n, m) = \prod p_i^{\min(a_i,\, b_i)}$ și
      $\text{lcm}(n, m) = \prod p_i^{\max(a_i,\, b_i)}$ \\
      Exemplu: $n = 2^3 \cdot 3^2,\ m = 2^1 \cdot 3^4 \;\Rightarrow\;
        \gcd = 2^1 \cdot 3^2 = 18,\;
        \text{lcm} = 2^3 \cdot 3^4 = 648$
  \end{itemize}
\end{frame}

\begin{frame}[fragile]{Algoritmul lui Euclid \textemdash{} implementare iterativă}
  \begin{block}{Explicația implementării}
    \begin{itemize}
      \item Se pornește de la $\gcd(a, b) = \gcd(a \bmod b,\ b) = \gcd(b,\ a \bmod b)$
      \item Cât timp $b > 0$, se calculează restul $r = a \bmod b$;
        apoi noul $a$ devine $b$, iar noul $b$ devine $r$
      \item După fiecare pas, noul $b$ este mai mic decât noul $a$ ($r < b$), deci primul
        termen rămâne cel mai mare și operația $a \bmod b$ se poate efectua în continuare
      \item Dacă inițial $a < b$, atunci $a \bmod b = a$, iar primul pas doar interschimbă
        valorile: $(a, b) \rightarrow (b, a)$
      \item Când $b = 0$, rezultatul este $a$ (cazul de bază $\gcd(a, 0) = a$)
    \end{itemize}
  \end{block}
  \begin{uncoverenv}<2->
    \begin{lstlisting}
int gcd(int a, int b) {
    while (b != 0) {
        int r = a % b;
        a = b;
        b = r;
    }
    return a;
}
    \end{lstlisting}
    \begin{itemize}
      \item Timp: $O(\log \min(a, b))$ \quad Spațiu: $O(1)$
    \end{itemize}
  \end{uncoverenv}
\end{frame}

\begin{frame}[fragile]{Cel mai mic multiplu comun \textemdash{} calcul}
  \begin{intrebare}
    Cum se calculează $\text{lcm}(a, b)$ dacă se cunoaște $\gcd(a, b)$?
  \end{intrebare}
  \uncover<2->{
    Din $a \cdot b = \gcd(a, b) \cdot \text{lcm}(a, b)$ rezultă
    \lstinline[keepspaces=true]|lcm = (a / gcd(a, b)) * b|.
  }
  \begin{alertblock}<3->{Ordinea operațiilor}
    Este indicat ca împărțirea să se facă prima: $a \cdot b$ poate produce overflow chiar dacă
    $\text{lcm}(a, b)$ încape în tipul de date, iar astfel rezultatul rămâne corect.
  \end{alertblock}
\end{frame}

\begin{frame}{Complexitatea algoritmului lui Euclid}
  \begin{intrebare}
    Cât devine $a$ după \textbf{doi} pași ai algoritmului?
  \end{intrebare}
  \begin{hint}<2->
    Se separă cazurile $b \leq a/2$ și $b > a/2$. Cât este $a \bmod b$ în fiecare caz?
  \end{hint}
  \begin{itemize}
    \item<3-> dacă $b \leq a/2$: $a \bmod b < b \leq a/2$
    \item<3-> dacă $b > a/2$: $a \bmod b = a - b < a/2$
  \end{itemize}
  \begin{alertblock}<4->{Concluzie}
    În ambele cazuri $a \bmod b < a/2$, deci după \textbf{doi} pași cel mai mare număr
    se înjumătățește cel puțin $\Rightarrow$ $O(\log \min(a, b))$ pași.
  \end{alertblock}
  \beginnote{$\log$ = numărul de înjumătățiri succesive până când un număr ajunge la $1$; $\log_2 10^{18} \approx 60$}
\end{frame}

% ================================================================
\section{Algoritmul lui Euclid extins}

\begin{frame}{Identitatea lui Bézout}
  \begin{block}{Identitatea lui Bézout}
    Pentru orice $a, b \in \mathbb{Z}$, există $x, y \in \mathbb{Z}$ astfel încât:
    \[ a \cdot x + b \cdot y = \gcd(a, b) \]
  \end{block}
  Exemplu: $22 \cdot 3 + 13 \cdot (-5) = 66 - 65 = 1 = \gcd(22, 13)$.
  \begin{intrebare}
    Cum se determină și numerele $x, y$, nu doar $\gcd$?
  \end{intrebare}
  \uncover<2->{
    Algoritmul lui Euclid reduce $(a, b)$ la o pereche mai mică, $(b,\ a \bmod b)$.
    Soluția pentru perechea mai mică poate fi folosită pentru perechea inițială?
  }
  \beginnote{Secțiunea conține un concept ușor mai dificil \textemdash{} poate fi reluată după parcurgerea restului lecției.}
\end{frame}

\begin{frame}{Euclid extins \textemdash{} ideea}
  Fie $g = \gcd(a, b)$ și ecuația $a \cdot x + b \cdot y = g$.
  Se scrie $a = q \cdot b + r$, unde $q = \lfloor a / b \rfloor$ și $r = a \bmod b$.
  \begin{intrebare}
    Ce formă are ecuația după înlocuirea lui $a$ și gruparea termenilor care conțin $b$?
  \end{intrebare}
  \uncover<2->{
    \[
      (\underbrace{q \cdot b + r}_{a}) \cdot x + b \cdot y = g
      \;\;\Rightarrow\;\;
      b \cdot (q \cdot x + y) + r \cdot x = g
    \]
  }
  \uncover<3->{
    Cu notațiile $x' = q \cdot x + y$ și $y' = x$, ecuația devine $b \cdot x' + r \cdot y' = g$.
  }
  \begin{itemize}
    \item<4-> Forma este aceeași, cu $(b,\ r)$ în loc de $(a,\ b)$ \textemdash{} exact un pas din algoritmul lui Euclid.
    \item<4-> Reducerea continuă până la cazul de bază $b = 0$: $g \cdot x + 0 \cdot y = g \;\Rightarrow\; x = 1,\ y = 0$
  \end{itemize}
  \begin{alertblock}<5->{Reconstrucție la întoarcere}
    Din $x' = q \cdot x + y$ și $y' = x$ rezultă $x = y'$, \quad $y = x' - q \cdot y'$.
  \end{alertblock}
\end{frame}

\begin{frame}{Exemplu pe hârtie: $22x + 13y = \gcd(22, 13)$}
  \textbf{Pasul 1: Se aplică algoritmul lui Euclid}
  \begin{itemize}
    \item $22 = 1 \cdot 13 + 9$ \hfill {\footnotesize\textcolor{gray}{$\Rightarrow 9 = 22 - 1 \cdot 13$}}
    \item $13 = 1 \cdot 9 + 4$ \hfill {\footnotesize\textcolor{gray}{$\Rightarrow 4 = 13 - 1 \cdot 9$}}
    \item $9 = 2 \cdot 4 + 1$ \hfill {\footnotesize\textcolor{gray}{$\Rightarrow 1 = 9 - 2 \cdot 4$}}
    \item $4 = 4 \cdot 1 + 0 \;\Rightarrow\; \gcd(22, 13) = 1$
  \end{itemize}
  \begin{intrebare}
    Cum se exprimă $1$ folosind doar $22$ și $13$?
  \end{intrebare}
  \uncover<2->{\textbf{Pasul 2: Substituție în ordine inversă, de jos în sus}}
  \begin{itemize}
    \item<2-> $1 = 1 \cdot \hla{9} - 2 \cdot \hlb{4}$
    \item<3-> $1 = 1 \cdot \hla{9} - 2 \cdot (\hlb{13} - 1 \cdot \hla{9}) = 3 \cdot \hla{9} - 2 \cdot \hlb{13}$
    \item<4-> $1 = 3 \cdot (1 \cdot \hla{22} - 1 \cdot \hlb{13}) - 2 \cdot \hlb{13} = 3 \cdot \hla{22} - 5 \cdot \hlb{13}$
  \end{itemize}
  \begin{alertblock}<5->{Soluție}
    $x = 3,\ y = -5$ \quad \checkmark \quad ($22 \cdot 3 + 13 \cdot (-5) = 66 - 65 = 1$)
  \end{alertblock}
\end{frame}

\begin{frame}[fragile]{Euclid extins \textemdash{} implementare iterativă}
  \begin{lstlisting}
// Gaseste x, y cu: a*x + b*y = gcd(a, b)
int gcd_ext(int a, int b, int &x, int &y) {
    int x0 = 1, y0 = 0;
    int x1 = 0, y1 = 1;
    while (b != 0) {
        int q  = a / b;
        int xn = x0 - q * x1;
        int yn = y0 - q * y1;
        int r  = a % b;
        a = b;  b = r;
        x0 = x1; y0 = y1;
        x1 = xn; y1 = yn;
    }
    x = x0; y = y0;
    return a;
}
  \end{lstlisting}
  \begin{itemize}
    \item \lstinline|int &x| = parametru transmis prin referință: modificările din funcție sunt vizibile în apelant
    \item Timp: $O(\log \min(a, b))$ \quad Spațiu: $O(1)$
  \end{itemize}
\end{frame}

\begin{frame}{Inițializarea coeficienților $x_0, y_0, x_1, y_1$}
  Fie $a_0, b_0$ valorile inițiale. La orice pas al buclei sunt menținute relațiile:
  \[
    a = a_0 \cdot x_0 + b_0 \cdot y_0 \qquad b = a_0 \cdot x_1 + b_0 \cdot y_1
  \]
  \begin{intrebare}
    Ce valori trebuie să aibă $x_0, y_0, x_1, y_1$ înainte de primul pas?
  \end{intrebare}
  \begin{itemize}
    \item<2-> Înainte de primul pas, $a = a_0$ și $b = b_0$.
    \item<2-> $a_0 = a_0 \cdot 1 + b_0 \cdot 0$ și $b_0 = a_0 \cdot 0 + b_0 \cdot 1$,
      deci $x_0 = 1,\ y_0 = 0,\ x_1 = 0,\ y_1 = 1$.
    \item<3-> La final $b = 0$ și $a = \gcd$, deci $x_0, y_0$ sunt coeficienții căutați.
  \end{itemize}
\end{frame}

% ================================================================
\section{Exponențiere rapidă}

\begin{frame}{Exponențiere rapidă \textemdash{} ideea}
  \begin{block}{Problema}
    Calculul eficient al lui $a^n$. Varianta naivă ($n$ înmulțiri) are complexitate $O(n)$.
  \end{block}
  \begin{intrebare}
    Cum se poate calcula $2^{10}$ cu mai puțin de $9$ înmulțiri?
  \end{intrebare}
  \begin{itemize}
    \item<2-> $2^{10} = (2^5)^2$: $2^5$ se calculează o singură dată și se ridică la pătrat
    \item<2-> Aceeași jumătate nu se mai calculează de două ori: exponentul se înjumătățește
  \end{itemize}
  \begin{intrebare}<3->
    Cum se procedează dacă exponentul este impar, de exemplu $2^5$?
  \end{intrebare}
  \begin{itemize}
    \item<4-> Se scoate o bază în afară: $2^5 = 2 \cdot 2^4$, iar $2^4$ are exponent par
  \end{itemize}
  \uncover<4->{
    \[
      a^n = \begin{cases}
        1 & n = 0 \\
        \left(a^{n/2}\right)^2 & n \text{ par} \\
        a \cdot a^{n-1} & n \text{ impar}
      \end{cases}
    \]
  }
\end{frame}

\begin{frame}{Numărul de pași}
  Exemplu: $2^{10} \rightarrow 2^5 \rightarrow 2^4 \rightarrow 2^2 \rightarrow 2^1 \rightarrow 2^0$
  \quad (5 pași în loc de 10)
  \begin{intrebare}
    Câți pași sunt necesari pentru $n = 10^9$?
  \end{intrebare}
  \begin{itemize}
    \item<2-> $10^9$ se împarte la $2$ de circa $30$ de ori, deoarece $2^{30} \approx 10^9$
    \item<2-> La fiecare înjumătățire se face cel mult o înmulțire suplimentară pentru exponentul impar,
      deci cel mult circa $2 \cdot 30$ înmulțiri
  \end{itemize}
  \begin{alertblock}<3->{Concluzie}
    Complexitate $O(\log n)$, în loc de $10^9$ pași.
  \end{alertblock}
\end{frame}

\begin{frame}[fragile]{Exponențiere rapidă \textemdash{} implementare iterativă}
  \begin{lstlisting}
long long power(long long base, int exp) {
    long long result = 1;
    while (exp > 0) {
        if (exp % 2 == 1)
            result *= base;
        base *= base;
        exp /= 2;
    }
    return result;
}
  \end{lstlisting}
  \begin{itemize}
    \item Timp: $O(\log n)$ \quad Spațiu: $O(1)$
    \item \alert{Atenție:} \lstinline|base * base| poate da overflow pentru baze mari \textemdash{}
      în concursuri se lucrează de obicei modulo un număr prim
  \end{itemize}
\end{frame}

\begin{frame}{Urmărirea execuției: \texttt{power(3, 5)}}
  \begin{intrebare}
    Ce valori au \texttt{result}, \texttt{base} și \texttt{exp} la finalul fiecărei iterații?
  \end{intrebare}
  \begin{center}
    \begin{tabular}{cccc}
      \toprule
      Iterația & \texttt{result} & \texttt{base} & \texttt{exp} \\
      \midrule
      început & $1$ & $3$ & $5$ \\
      \at{2}{$1$ (\texttt{exp} $= 5$, impar)} & \at{2}{$3$} & \at{2}{$9$} & \at{2}{$2$} \\
      \at{3}{$2$ (\texttt{exp} $= 2$, par)}   & \at{3}{$3$} & \at{3}{$81$} & \at{3}{$1$} \\
      \at{4}{$3$ (\texttt{exp} $= 1$, impar)} & \at{4}{$243$} & \at{4}{$6561$} & \at{4}{$0$} \\
      \bottomrule
    \end{tabular}
  \end{center}
  \begin{alertblock}<5->{Rezultat}
    \texttt{power(3, 5)} $= 243 = 3^5$ \quad \checkmark
  \end{alertblock}
\end{frame}

\begin{frame}{Legătura cu scrierea binară}
  În exemplul \texttt{power(3, 5)}: $5 = 101_{(2)}$, iar \texttt{base} a luat valorile $3^1,\ 3^2,\ 3^4$.
  \begin{intrebare}
    Ce legătură există între algoritm și scrierea lui $n$ în baza $2$?
  \end{intrebare}
  \begin{hint}<2->
    Se descompune $5 = 4 + 1$ și se scrie $3^5$ ca produs de puteri ale lui $3$.
  \end{hint}
  \begin{itemize}
    \item<3-> $3^5 = 3^4 \cdot 3^1$; \texttt{base} parcurge $a^{2^0}, a^{2^1}, a^{2^2}, \ldots$, iar
      \texttt{result} se înmulțește doar cu bazele corespunzătoare \textbf{biților egali cu $1$}
      din scrierea binară a lui $n$
  \end{itemize}
  \begin{alertblock}<4->{Concluzie}
    Numărul de biți ai lui $n$ este $\approx \log_2 n$, deci $O(\log n)$ iterații.
  \end{alertblock}
\end{frame}

% ================================================================
\section{Bonus \textemdash{} Implementări recursive}

\begin{frame}{Recursivitatea}
  \begin{block}{Definiție}
    O funcție este \textbf{recursivă} dacă se apelează pe ea însăși.
  \end{block}
  Orice funcție recursivă are nevoie de:
  \begin{itemize}
    \item Un \textbf{caz de bază} \textemdash{} condiția de oprire
    \item Un \textbf{apel recursiv} pe o subproblemă mai mică
  \end{itemize}
  Exemplu: $n! = n \cdot (n-1)!$, cu cazul de bază $0! = 1$
  \begin{intrebare}
    Ce se întâmplă dacă lipsește cazul de bază?
  \end{intrebare}
  \begin{itemize}
    \item<2-> Funcția s-ar apela la nesfârșit până se umple stiva de apeluri \textemdash{}
      \textbf{stack overflow}
  \end{itemize}
  \begin{block}<3->{\textcolor{advAlert}{Atenție} \textemdash{} dezavantajele recursivității}
    Overhead de apeluri de funcție (ușor mai lent) și risc de \textbf{stack overflow} pentru intrări
    foarte mari. Codul recursiv reflectă direct definiția matematică, dar în concursuri
    este preferată varianta iterativă.
  \end{block}
\end{frame}

\begin{frame}[fragile]{Euclid \textemdash{} implementare recursivă}
  \begin{intrebare}
    Cum se scrie $\gcd$ recursiv pornind de la $\gcd(a, b) = \gcd(b,\ a \bmod b)$?
    Care este cazul de bază?
  \end{intrebare}
  \begin{uncoverenv}<2->
    \begin{lstlisting}
int gcd(int a, int b) {
    if (b == 0)
        return a;
    return gcd(b, a % b);
}
    \end{lstlisting}
    \begin{itemize}
      \item Cazul de bază este $\gcd(a, 0) = a$
      \item Codul reflectă direct proprietatea $\gcd(a, b) = \gcd(b,\ a \bmod b)$
      \item Aceeași complexitate ca varianta iterativă: $O(\log \min(a, b))$
    \end{itemize}
  \end{uncoverenv}
\end{frame}

\begin{frame}[fragile]{Euclid extins \textemdash{} implementare recursivă}
  \begin{intrebare}
    Ce valori au $x$ și $y$ în cazul de bază $b = 0$? Cum se reconstruiesc la întoarcere?
  \end{intrebare}
  \begin{itemize}
    \item<2-> Cazul de bază: $a \cdot x + 0 \cdot y = a$, deci $x = 1,\ y = 0$
    \item<2-> La întoarcere: $x = y'$, \quad $y = x' - q \cdot y'$
  \end{itemize}
  \begin{uncoverenv}<3->
    \begin{lstlisting}
int gcd_ext(int a, int b, int &x, int &y) {
    if (b == 0) {
        x = 1; y = 0;
        return a;
    }
    int x1, y1;
    int d = gcd_ext(b, a % b, x1, y1);
    x = y1;
    y = x1 - (a / b) * y1;
    return d;
}
    \end{lstlisting}
    \begin{itemize}
      \item Timp: $O(\log \min(a, b))$ \quad Spațiu: $O(\log \min(a, b))$ \textemdash{} stiva de apeluri
    \end{itemize}
  \end{uncoverenv}
\end{frame}

\begin{frame}[fragile]{Exponențiere rapidă \textemdash{} implementare recursivă}
  \begin{lstlisting}
long long power(long long base, int exp) {
    if (exp == 0)
        return 1;
    if (exp % 2 == 0) {
        long long half = power(base, exp / 2);
        return half * half;
    }
    return base * power(base, exp - 1);
}
  \end{lstlisting}
  \begin{itemize}
    \item Timp: $O(\log n)$ \quad Spațiu: $O(\log n)$ \textemdash{} stiva de apeluri
  \end{itemize}
  \begin{intrebare}
    De ce nu se scrie direct
    \lstinline[keepspaces=true]|return power(base, exp / 2) * power(base, exp / 2);|?
  \end{intrebare}
  \begin{itemize}
    \item<2-> Două apeluri la fiecare nivel $\Rightarrow$ numărul de apeluri se dublează la fiecare nivel
    \item<2-> În total $O(n)$ apeluri \textemdash{} câștigul se pierde; cu \lstinline|half * half|
      se face \textbf{un singur} apel pe nivel: $O(\log n)$
  \end{itemize}
\end{frame}

% ================================================================
\section{Recapitulare}

\begin{frame}{Recapitulare \textemdash{} Complexități}
  \begin{center}
  \begin{tabular}{lll}
    \toprule
    Algoritm & Timp & Spațiu \\
    \midrule
    $\gcd$ iterativ                & $O(\log \min(a,b))$ & $O(1)$ \\
    $\gcd$ extins iterativ         & $O(\log \min(a,b))$ & $O(1)$ \\
    $\gcd$ recursiv                & $O(\log \min(a,b))$ & $O(\log \min(a,b))$ \\
    $\gcd$ extins recursiv         & $O(\log \min(a,b))$ & $O(\log \min(a,b))$ \\
    Exponențiere rapidă iterativă  & $O(\log n)$         & $O(1)$ \\
    Exponențiere rapidă recursivă  & $O(\log n)$         & $O(\log n)$ \\
    \bottomrule
  \end{tabular}
  \end{center}
\end{frame}

\begin{frame}{Probleme de antrenament}
  \begin{block}{AlgoCode}
    \begin{enumerate}
      \item \href{https://code.algolymp.com/problems/1557}{\textbf{GCD and LCM}}
            \textcolor{gray}{(AlgoCode \#1557)}
      \item \href{https://code.algolymp.com/problems/4230}{\textbf{GCD and LCM Queries}}
            \textcolor{gray}{(AlgoCode \#4230)}
      \item \href{https://code.algolymp.com/problems/4270}{\textbf{GCD Power Queries}}
            \textcolor{gray}{(AlgoCode \#4270)}
    \end{enumerate}
  \end{block}
  \begin{block}{Altele}
    \begin{enumerate}
      \item \href{https://www.pbinfo.ro/probleme/58/cmmdc}{\textbf{cmmdc}}
            \textcolor{gray}{(pbinfo.ro)}
      \item \href{https://www.pbinfo.ro/probleme/59/cmmmc}{\textbf{cmmmc}}
            \textcolor{gray}{(pbinfo.ro)}
      \item \href{https://www.infoarena.ro/problema/euclid3}{\textbf{euclid3}}
            \textcolor{gray}{(infoarena.ro)}
      \item \href{https://infoarena.ro/problema/lgput}{\textbf{lgput}}
            \textcolor{gray}{(infoarena.ro)}
    \end{enumerate}
  \end{block}
\end{frame}

\begin{frame}{Resurse}
  \begin{enumerate}
    \item \href{https://www.infoarena.ro/algoritmul-lui-euclid}{\textbf{infoarena.ro}}
          \textemdash{} Algoritmul lui Euclid (RO)
    \item \href{https://cp-algorithms.com/algebra/euclid-algorithm.html}{\textbf{cp-algorithms.com}}
          \textemdash{} Euclidean algorithm (EN)
    \item \href{https://cp-algorithms.com/algebra/extended-euclid-algorithm.html}{\textbf{cp-algorithms.com}}
          \textemdash{} Extended Euclidean algorithm (EN)
    \item \href{https://cp-algorithms.com/algebra/binary-exp.html}{\textbf{cp-algorithms.com}}
          \textemdash{} Binary exponentiation (EN)
    \item \href{https://www.pbinfo.ro/articole/3873/recursivitate}{\textbf{pbinfo.ro}}
          \textemdash{} Recursivitate (RO)
    \item \href{https://cppreference.com}{\textbf{cppreference.com}}
          \textemdash{} documentație completă C++
  \end{enumerate}
\end{frame}

\end{document}
"
% ================================================================
\documentclass[9pt]{beamer}
\usepackage{../../../resources/themes/algolymp-dense}

% Evidențiere a doi termeni care se urmăresc prin substituții
\newcommand{\hla}[1]{\textbf{\textcolor{advAlert}{#1}}}
\newcommand{\hlb}[1]{\textbf{\textcolor{advAccent}{#1}}}

% Mod interactiv: întrebarea apare prima, apoi soluția se dezvăluie pe bucăți
% (\item<2->, \uncover<3->{...}, \begin{alertblock}<4->{...}); hint-ul se
% folosește doar unde elevii chiar trebuie împinși să se gândească.
% Totul se afișează cumulativ: ce apare la pasul 2 rămâne vizibil după.
\newenvironment<>{intrebare}{%
  \begin{uncoverenv}#1%
  \begin{block}{Întrebare}}{\end{block}\end{uncoverenv}}
\newenvironment<>{hint}{%
  \begin{uncoverenv}#1%
  \setbeamercolor{block title}{bg=advAlert!15,fg=advAlert}%
  \setbeamercolor{block body}{bg=advAlert!5,fg=black}%
  \begin{block}{Hint}}{\end{block}\end{uncoverenv}}

% Celulă de tabel vizibilă de la pasul #1
\newcommand{\at}[2]{\uncover<#1->{#2}}

\title{Algoritmul lui Euclid \textbar\ Exponențiere rapidă}
\subtitle{}
\author{Andrei Mihai Ciobanu}
\institute{Algolymp Summer School}
\date{}
\titlegraphic{\includegraphics[width=3cm]{../../../resources/algolymp-logo.png}}
\setbeamertemplate{frame footer}{Algolymp Summer School}

% ================================================================
\begin{document}

\begin{frame}\titlepage\end{frame}

\begin{frame}{Cuprins}\tableofcontents\end{frame}

% ================================================================
\section{Algoritmul lui Euclid}

\begin{frame}{Cel mai mare divizor comun}
  \begin{block}{Cel mai mare divizor comun}
    $\text{cmmdc}(a, b)$ = cel mai mare număr natural $d$ astfel încât $d \mid a$ și $d \mid b$.
    Se mai notează și $\gcd(a, b)$ (\textit{greatest common divisor}) sau doar $(a, b)$.
  \end{block}
  \begin{block}{Cel mai mic multiplu comun}
    $\text{cmmmc}(a, b)$ = cel mai mic număr natural $m$ astfel încât $a \mid m$ și $b \mid m$.
    Se mai notează și $\text{lcm}(a, b)$ (\textit{lowest common multiple}) sau doar $[a, b]$.
  \end{block}
  \smallskip
  Exemple:
  \begin{itemize}
    \item $\gcd(12, 8) = 4$ \quad $\gcd(17, 5) = 1$ \quad $\gcd(100, 75) = 25$
    \item $\text{lcm}(4, 6) = 12$ \quad $\text{lcm}(3, 5) = 15$
  \end{itemize}
\end{frame}

\begin{frame}{Calculul cmmdc}
  \begin{intrebare}
    Cum se calculează $\gcd(84, 36)$? Dar $\gcd(10^{18},\ 6 \cdot 10^{17})$?
  \end{intrebare}
  \begin{itemize}
    \item<2-> Se listează divizorii lui $a$ în $O(\sqrt{a})$ și se testează fiecare pe $b$.
      Pentru $\gcd(84, 36) = 12$ metoda funcționează.
    \item<2-> Pentru numere de ordinul $10^{18}$ rezultă circa $10^9$ operații \textemdash{} prea lent.
  \end{itemize}
  \begin{alertblock}<3->{Necesitatea unei metode mai rapide}
    \textbf{Algoritmul lui Euclid}: sub $100$ de pași pentru numere de $10^{18}$.
  \end{alertblock}
\end{frame}

\begin{frame}{Ideea de bază: $\gcd(a, b) = \gcd(a - b,\ b)$}
  \begin{intrebare}
    Fie $d$ un divizor comun al lui $a$ și $b$. Ce se poate afirma despre $a - b$?
  \end{intrebare}
  \begin{hint}<2->
    Se scriu $a = d \cdot x$ și $b = d \cdot y$, apoi se calculează $a - b$.
  \end{hint}
  \begin{itemize}
    \item<3-> $a - b = d \cdot (x - y)$, deci $d \mid a - b$.
    \item<3-> Reciproc, dacă $d \mid a - b$ și $d \mid b$, atunci $d \mid (a - b) + b = a$.
  \end{itemize}
  \begin{alertblock}<4->{Concluzie}
    Perechile $(a, b)$ și $(a - b,\ b)$ au \textbf{aceiași divizori comuni}, deci
    $\gcd(a, b) = \gcd(a - b,\ b)$.
  \end{alertblock}
\end{frame}

\begin{frame}{Scăderi repetate}
  Regula $\gcd(a, b) = \gcd(a - b, b)$ se aplică repetat până la cazul de bază $\gcd(a, 0) = a$.
  \begin{intrebare}
    Cât este $\gcd(12, 8)$ dacă se folosesc \textbf{doar} scăderi?
  \end{intrebare}
  \uncover<2->{
    \[
      \gcd(12, 8) = \gcd(4, 8) = \gcd(4, 4) = \gcd(0, 4) = 4 \quad \checkmark
    \]
  }
  \begin{intrebare}<3->
    Câți pași sunt necesari pentru $\gcd(1000, 1)$?
  \end{intrebare}
  \uncover<4->{
    \[
      \gcd(1000, 1) \rightarrow \gcd(999, 1) \rightarrow \gcd(998, 1) \rightarrow \cdots
      \rightarrow \gcd(0, 1) = 1
    \]
    \textbf{1000 de pași} pentru un rezultat evident \textemdash{} procedeul este prea lent pentru valori mari.
  }
\end{frame}

\begin{frame}{De la scăderi la rest}
  \begin{intrebare}
    Cum pot fi înlocuite scăderile repetate ale lui $b$ cu o singură operație?
  \end{intrebare}
  \uncover<2->{
    Prin împărțirea cu rest, $a = q \cdot b + r$, cu $0 \leq r < b$. Restul $r$ se notează
    $a \bmod b$ (restul împărțirii lui $a$ la $b$):
    \[
      \gcd(a, b) = \gcd(a - b, b) = \gcd(a - 2b, b) = \cdots = \gcd(a \bmod b,\ b)
    \]
  }
  \begin{alertblock}<3->{Algoritmul lui Euclid}
    $\gcd(a, b) = \gcd(b,\ a \bmod b)$, \quad cu cazul de bază $\gcd(a, 0) = a$.
  \end{alertblock}
  \uncover<4->{Exemplu: $\gcd(48, 18)$}
  \begin{itemize}
    \item<4-> $\gcd(48, 18) = \gcd(18, 12)$ \quad ($48 \bmod 18 = 12$)
    \item<4-> $\gcd(18, 12) = \gcd(12, 6)$ \quad\ ($18 \bmod 12 = 6$)
    \item<5-> $\gcd(12,\ 6) = \gcd(6,\ 0)$ \quad\ \ ($12 \bmod 6 = 0$)
    \item<5-> $\gcd(6, 0) = 6$ \checkmark
  \end{itemize}
\end{frame}

\begin{frame}{Proprietăți}
  \begin{itemize}
    \item $\gcd(a, 0) = \gcd(a, a) = |a|$ \quad (\textit{cazuri de bază})
    \item $\gcd(a, b) = \gcd(b, a)$ \quad (\textit{comutativitate})
    \item $\gcd(a, \gcd(b, c)) = \gcd(\gcd(a, b), c)$ \quad (\textit{asociativitate})
    \item $\gcd(k \cdot a, k \cdot b) = k \cdot \gcd(a, b)$, pentru $k > 0$
    \item dacă $\gcd(a, b) = d$, atunci $\gcd\!\left(\dfrac{a}{d},\ \dfrac{b}{d}\right) = 1$
    \item $a \cdot b = \gcd(a, b) \cdot \text{lcm}(a, b)$
    \item dacă $\gcd(a, b) = 1$, se spune că $a$ și $b$ sunt \textit{prime între ele} (\textit{coprime})
    \item $\gcd(a, b) = \gcd(a \pm k \cdot b,\ b)$, $\forall\, k \in \mathbb{N}$ \quad
      $\Rightarrow$ \quad $\gcd(a, b) = \gcd(b,\ a \bmod b)$
    \item \textbf{Prin descompunere în factori primi:} dacă
      $n = \prod p_i^{a_i}$ și $m = \prod p_i^{b_i}$, atunci
      $\gcd(n, m) = \prod p_i^{\min(a_i,\, b_i)}$ și
      $\text{lcm}(n, m) = \prod p_i^{\max(a_i,\, b_i)}$ \\
      Exemplu: $n = 2^3 \cdot 3^2,\ m = 2^1 \cdot 3^4 \;\Rightarrow\;
        \gcd = 2^1 \cdot 3^2 = 18,\;
        \text{lcm} = 2^3 \cdot 3^4 = 648$
  \end{itemize}
\end{frame}

\begin{frame}[fragile]{Algoritmul lui Euclid \textemdash{} implementare iterativă}
  \begin{block}{Explicația implementării}
    \begin{itemize}
      \item Se pornește de la $\gcd(a, b) = \gcd(a \bmod b,\ b) = \gcd(b,\ a \bmod b)$
      \item Cât timp $b > 0$, se calculează restul $r = a \bmod b$;
        apoi noul $a$ devine $b$, iar noul $b$ devine $r$
      \item După fiecare pas, noul $b$ este mai mic decât noul $a$ ($r < b$), deci primul
        termen rămâne cel mai mare și operația $a \bmod b$ se poate efectua în continuare
      \item Dacă inițial $a < b$, atunci $a \bmod b = a$, iar primul pas doar interschimbă
        valorile: $(a, b) \rightarrow (b, a)$
      \item Când $b = 0$, rezultatul este $a$ (cazul de bază $\gcd(a, 0) = a$)
    \end{itemize}
  \end{block}
  \begin{uncoverenv}<2->
    \begin{lstlisting}
int gcd(int a, int b) {
    while (b != 0) {
        int r = a % b;
        a = b;
        b = r;
    }
    return a;
}
    \end{lstlisting}
    \begin{itemize}
      \item Timp: $O(\log \min(a, b))$ \quad Spațiu: $O(1)$
    \end{itemize}
  \end{uncoverenv}
\end{frame}

\begin{frame}[fragile]{Cel mai mic multiplu comun \textemdash{} calcul}
  \begin{intrebare}
    Cum se calculează $\text{lcm}(a, b)$ dacă se cunoaște $\gcd(a, b)$?
  \end{intrebare}
  \uncover<2->{
    Din $a \cdot b = \gcd(a, b) \cdot \text{lcm}(a, b)$ rezultă
    \lstinline[keepspaces=true]|lcm = (a / gcd(a, b)) * b|.
  }
  \begin{alertblock}<3->{Ordinea operațiilor}
    Este indicat ca împărțirea să se facă prima: $a \cdot b$ poate produce overflow chiar dacă
    $\text{lcm}(a, b)$ încape în tipul de date, iar astfel rezultatul rămâne corect.
  \end{alertblock}
\end{frame}

\begin{frame}{Complexitatea algoritmului lui Euclid}
  \begin{intrebare}
    Cât devine $a$ după \textbf{doi} pași ai algoritmului?
  \end{intrebare}
  \begin{hint}<2->
    Se separă cazurile $b \leq a/2$ și $b > a/2$. Cât este $a \bmod b$ în fiecare caz?
  \end{hint}
  \begin{itemize}
    \item<3-> dacă $b \leq a/2$: $a \bmod b < b \leq a/2$
    \item<3-> dacă $b > a/2$: $a \bmod b = a - b < a/2$
  \end{itemize}
  \begin{alertblock}<4->{Concluzie}
    În ambele cazuri $a \bmod b < a/2$, deci după \textbf{doi} pași cel mai mare număr
    se înjumătățește cel puțin $\Rightarrow$ $O(\log \min(a, b))$ pași.
  \end{alertblock}
  \beginnote{$\log$ = numărul de înjumătățiri succesive până când un număr ajunge la $1$; $\log_2 10^{18} \approx 60$}
\end{frame}

% ================================================================
\section{Algoritmul lui Euclid extins}

\begin{frame}{Identitatea lui Bézout}
  \begin{block}{Identitatea lui Bézout}
    Pentru orice $a, b \in \mathbb{Z}$, există $x, y \in \mathbb{Z}$ astfel încât:
    \[ a \cdot x + b \cdot y = \gcd(a, b) \]
  \end{block}
  Exemplu: $22 \cdot 3 + 13 \cdot (-5) = 66 - 65 = 1 = \gcd(22, 13)$.
  \begin{intrebare}
    Cum se determină și numerele $x, y$, nu doar $\gcd$?
  \end{intrebare}
  \uncover<2->{
    Algoritmul lui Euclid reduce $(a, b)$ la o pereche mai mică, $(b,\ a \bmod b)$.
    Soluția pentru perechea mai mică poate fi folosită pentru perechea inițială?
  }
  \beginnote{Secțiunea conține un concept ușor mai dificil \textemdash{} poate fi reluată după parcurgerea restului lecției.}
\end{frame}

\begin{frame}{Euclid extins \textemdash{} ideea}
  Fie $g = \gcd(a, b)$ și ecuația $a \cdot x + b \cdot y = g$.
  Se scrie $a = q \cdot b + r$, unde $q = \lfloor a / b \rfloor$ și $r = a \bmod b$.
  \begin{intrebare}
    Ce formă are ecuația după înlocuirea lui $a$ și gruparea termenilor care conțin $b$?
  \end{intrebare}
  \uncover<2->{
    \[
      (\underbrace{q \cdot b + r}_{a}) \cdot x + b \cdot y = g
      \;\;\Rightarrow\;\;
      b \cdot (q \cdot x + y) + r \cdot x = g
    \]
  }
  \uncover<3->{
    Cu notațiile $x' = q \cdot x + y$ și $y' = x$, ecuația devine $b \cdot x' + r \cdot y' = g$.
  }
  \begin{itemize}
    \item<4-> Forma este aceeași, cu $(b,\ r)$ în loc de $(a,\ b)$ \textemdash{} exact un pas din algoritmul lui Euclid.
    \item<4-> Reducerea continuă până la cazul de bază $b = 0$: $g \cdot x + 0 \cdot y = g \;\Rightarrow\; x = 1,\ y = 0$
  \end{itemize}
  \begin{alertblock}<5->{Reconstrucție la întoarcere}
    Din $x' = q \cdot x + y$ și $y' = x$ rezultă $x = y'$, \quad $y = x' - q \cdot y'$.
  \end{alertblock}
\end{frame}

\begin{frame}{Exemplu pe hârtie: $22x + 13y = \gcd(22, 13)$}
  \textbf{Pasul 1: Se aplică algoritmul lui Euclid}
  \begin{itemize}
    \item $22 = 1 \cdot 13 + 9$ \hfill {\footnotesize\textcolor{gray}{$\Rightarrow 9 = 22 - 1 \cdot 13$}}
    \item $13 = 1 \cdot 9 + 4$ \hfill {\footnotesize\textcolor{gray}{$\Rightarrow 4 = 13 - 1 \cdot 9$}}
    \item $9 = 2 \cdot 4 + 1$ \hfill {\footnotesize\textcolor{gray}{$\Rightarrow 1 = 9 - 2 \cdot 4$}}
    \item $4 = 4 \cdot 1 + 0 \;\Rightarrow\; \gcd(22, 13) = 1$
  \end{itemize}
  \begin{intrebare}
    Cum se exprimă $1$ folosind doar $22$ și $13$?
  \end{intrebare}
  \uncover<2->{\textbf{Pasul 2: Substituție în ordine inversă, de jos în sus}}
  \begin{itemize}
    \item<2-> $1 = 1 \cdot \hla{9} - 2 \cdot \hlb{4}$
    \item<3-> $1 = 1 \cdot \hla{9} - 2 \cdot (\hlb{13} - 1 \cdot \hla{9}) = 3 \cdot \hla{9} - 2 \cdot \hlb{13}$
    \item<4-> $1 = 3 \cdot (1 \cdot \hla{22} - 1 \cdot \hlb{13}) - 2 \cdot \hlb{13} = 3 \cdot \hla{22} - 5 \cdot \hlb{13}$
  \end{itemize}
  \begin{alertblock}<5->{Soluție}
    $x = 3,\ y = -5$ \quad \checkmark \quad ($22 \cdot 3 + 13 \cdot (-5) = 66 - 65 = 1$)
  \end{alertblock}
\end{frame}

\begin{frame}[fragile]{Euclid extins \textemdash{} implementare iterativă}
  \begin{lstlisting}
// Gaseste x, y cu: a*x + b*y = gcd(a, b)
int gcd_ext(int a, int b, int &x, int &y) {
    int x0 = 1, y0 = 0;
    int x1 = 0, y1 = 1;
    while (b != 0) {
        int q  = a / b;
        int xn = x0 - q * x1;
        int yn = y0 - q * y1;
        int r  = a % b;
        a = b;  b = r;
        x0 = x1; y0 = y1;
        x1 = xn; y1 = yn;
    }
    x = x0; y = y0;
    return a;
}
  \end{lstlisting}
  \begin{itemize}
    \item \lstinline|int &x| = parametru transmis prin referință: modificările din funcție sunt vizibile în apelant
    \item Timp: $O(\log \min(a, b))$ \quad Spațiu: $O(1)$
  \end{itemize}
\end{frame}

\begin{frame}{Inițializarea coeficienților $x_0, y_0, x_1, y_1$}
  Fie $a_0, b_0$ valorile inițiale. La orice pas al buclei sunt menținute relațiile:
  \[
    a = a_0 \cdot x_0 + b_0 \cdot y_0 \qquad b = a_0 \cdot x_1 + b_0 \cdot y_1
  \]
  \begin{intrebare}
    Ce valori trebuie să aibă $x_0, y_0, x_1, y_1$ înainte de primul pas?
  \end{intrebare}
  \begin{itemize}
    \item<2-> Înainte de primul pas, $a = a_0$ și $b = b_0$.
    \item<2-> $a_0 = a_0 \cdot 1 + b_0 \cdot 0$ și $b_0 = a_0 \cdot 0 + b_0 \cdot 1$,
      deci $x_0 = 1,\ y_0 = 0,\ x_1 = 0,\ y_1 = 1$.
    \item<3-> La final $b = 0$ și $a = \gcd$, deci $x_0, y_0$ sunt coeficienții căutați.
  \end{itemize}
\end{frame}

% ================================================================
\section{Exponențiere rapidă}

\begin{frame}{Exponențiere rapidă \textemdash{} ideea}
  \begin{block}{Problema}
    Calculul eficient al lui $a^n$. Varianta naivă ($n$ înmulțiri) are complexitate $O(n)$.
  \end{block}
  \begin{intrebare}
    Cum se poate calcula $2^{10}$ cu mai puțin de $9$ înmulțiri?
  \end{intrebare}
  \begin{itemize}
    \item<2-> $2^{10} = (2^5)^2$: $2^5$ se calculează o singură dată și se ridică la pătrat
    \item<2-> Aceeași jumătate nu se mai calculează de două ori: exponentul se înjumătățește
  \end{itemize}
  \begin{intrebare}<3->
    Cum se procedează dacă exponentul este impar, de exemplu $2^5$?
  \end{intrebare}
  \begin{itemize}
    \item<4-> Se scoate o bază în afară: $2^5 = 2 \cdot 2^4$, iar $2^4$ are exponent par
  \end{itemize}
  \uncover<4->{
    \[
      a^n = \begin{cases}
        1 & n = 0 \\
        \left(a^{n/2}\right)^2 & n \text{ par} \\
        a \cdot a^{n-1} & n \text{ impar}
      \end{cases}
    \]
  }
\end{frame}

\begin{frame}{Numărul de pași}
  Exemplu: $2^{10} \rightarrow 2^5 \rightarrow 2^4 \rightarrow 2^2 \rightarrow 2^1 \rightarrow 2^0$
  \quad (5 pași în loc de 10)
  \begin{intrebare}
    Câți pași sunt necesari pentru $n = 10^9$?
  \end{intrebare}
  \begin{itemize}
    \item<2-> $10^9$ se împarte la $2$ de circa $30$ de ori, deoarece $2^{30} \approx 10^9$
    \item<2-> La fiecare înjumătățire se face cel mult o înmulțire suplimentară pentru exponentul impar,
      deci cel mult circa $2 \cdot 30$ înmulțiri
  \end{itemize}
  \begin{alertblock}<3->{Concluzie}
    Complexitate $O(\log n)$, în loc de $10^9$ pași.
  \end{alertblock}
\end{frame}

\begin{frame}[fragile]{Exponențiere rapidă \textemdash{} implementare iterativă}
  \begin{lstlisting}
long long power(long long base, int exp) {
    long long result = 1;
    while (exp > 0) {
        if (exp % 2 == 1)
            result *= base;
        base *= base;
        exp /= 2;
    }
    return result;
}
  \end{lstlisting}
  \begin{itemize}
    \item Timp: $O(\log n)$ \quad Spațiu: $O(1)$
    \item \alert{Atenție:} \lstinline|base * base| poate da overflow pentru baze mari \textemdash{}
      în concursuri se lucrează de obicei modulo un număr prim
  \end{itemize}
\end{frame}

\begin{frame}{Urmărirea execuției: \texttt{power(3, 5)}}
  \begin{intrebare}
    Ce valori au \texttt{result}, \texttt{base} și \texttt{exp} la finalul fiecărei iterații?
  \end{intrebare}
  \begin{center}
    \begin{tabular}{cccc}
      \toprule
      Iterația & \texttt{result} & \texttt{base} & \texttt{exp} \\
      \midrule
      început & $1$ & $3$ & $5$ \\
      \at{2}{$1$ (\texttt{exp} $= 5$, impar)} & \at{2}{$3$} & \at{2}{$9$} & \at{2}{$2$} \\
      \at{3}{$2$ (\texttt{exp} $= 2$, par)}   & \at{3}{$3$} & \at{3}{$81$} & \at{3}{$1$} \\
      \at{4}{$3$ (\texttt{exp} $= 1$, impar)} & \at{4}{$243$} & \at{4}{$6561$} & \at{4}{$0$} \\
      \bottomrule
    \end{tabular}
  \end{center}
  \begin{alertblock}<5->{Rezultat}
    \texttt{power(3, 5)} $= 243 = 3^5$ \quad \checkmark
  \end{alertblock}
\end{frame}

\begin{frame}{Legătura cu scrierea binară}
  În exemplul \texttt{power(3, 5)}: $5 = 101_{(2)}$, iar \texttt{base} a luat valorile $3^1,\ 3^2,\ 3^4$.
  \begin{intrebare}
    Ce legătură există între algoritm și scrierea lui $n$ în baza $2$?
  \end{intrebare}
  \begin{hint}<2->
    Se descompune $5 = 4 + 1$ și se scrie $3^5$ ca produs de puteri ale lui $3$.
  \end{hint}
  \begin{itemize}
    \item<3-> $3^5 = 3^4 \cdot 3^1$; \texttt{base} parcurge $a^{2^0}, a^{2^1}, a^{2^2}, \ldots$, iar
      \texttt{result} se înmulțește doar cu bazele corespunzătoare \textbf{biților egali cu $1$}
      din scrierea binară a lui $n$
  \end{itemize}
  \begin{alertblock}<4->{Concluzie}
    Numărul de biți ai lui $n$ este $\approx \log_2 n$, deci $O(\log n)$ iterații.
  \end{alertblock}
\end{frame}

% ================================================================
\section{Bonus \textemdash{} Implementări recursive}

\begin{frame}{Recursivitatea}
  \begin{block}{Definiție}
    O funcție este \textbf{recursivă} dacă se apelează pe ea însăși.
  \end{block}
  Orice funcție recursivă are nevoie de:
  \begin{itemize}
    \item Un \textbf{caz de bază} \textemdash{} condiția de oprire
    \item Un \textbf{apel recursiv} pe o subproblemă mai mică
  \end{itemize}
  Exemplu: $n! = n \cdot (n-1)!$, cu cazul de bază $0! = 1$
  \begin{intrebare}
    Ce se întâmplă dacă lipsește cazul de bază?
  \end{intrebare}
  \begin{itemize}
    \item<2-> Funcția s-ar apela la nesfârșit până se umple stiva de apeluri \textemdash{}
      \textbf{stack overflow}
  \end{itemize}
  \begin{block}<3->{\textcolor{advAlert}{Atenție} \textemdash{} dezavantajele recursivității}
    Overhead de apeluri de funcție (ușor mai lent) și risc de \textbf{stack overflow} pentru intrări
    foarte mari. Codul recursiv reflectă direct definiția matematică, dar în concursuri
    este preferată varianta iterativă.
  \end{block}
\end{frame}

\begin{frame}[fragile]{Euclid \textemdash{} implementare recursivă}
  \begin{intrebare}
    Cum se scrie $\gcd$ recursiv pornind de la $\gcd(a, b) = \gcd(b,\ a \bmod b)$?
    Care este cazul de bază?
  \end{intrebare}
  \begin{uncoverenv}<2->
    \begin{lstlisting}
int gcd(int a, int b) {
    if (b == 0)
        return a;
    return gcd(b, a % b);
}
    \end{lstlisting}
    \begin{itemize}
      \item Cazul de bază este $\gcd(a, 0) = a$
      \item Codul reflectă direct proprietatea $\gcd(a, b) = \gcd(b,\ a \bmod b)$
      \item Aceeași complexitate ca varianta iterativă: $O(\log \min(a, b))$
    \end{itemize}
  \end{uncoverenv}
\end{frame}

\begin{frame}[fragile]{Euclid extins \textemdash{} implementare recursivă}
  \begin{intrebare}
    Ce valori au $x$ și $y$ în cazul de bază $b = 0$? Cum se reconstruiesc la întoarcere?
  \end{intrebare}
  \begin{itemize}
    \item<2-> Cazul de bază: $a \cdot x + 0 \cdot y = a$, deci $x = 1,\ y = 0$
    \item<2-> La întoarcere: $x = y'$, \quad $y = x' - q \cdot y'$
  \end{itemize}
  \begin{uncoverenv}<3->
    \begin{lstlisting}
int gcd_ext(int a, int b, int &x, int &y) {
    if (b == 0) {
        x = 1; y = 0;
        return a;
    }
    int x1, y1;
    int d = gcd_ext(b, a % b, x1, y1);
    x = y1;
    y = x1 - (a / b) * y1;
    return d;
}
    \end{lstlisting}
    \begin{itemize}
      \item Timp: $O(\log \min(a, b))$ \quad Spațiu: $O(\log \min(a, b))$ \textemdash{} stiva de apeluri
    \end{itemize}
  \end{uncoverenv}
\end{frame}

\begin{frame}[fragile]{Exponențiere rapidă \textemdash{} implementare recursivă}
  \begin{lstlisting}
long long power(long long base, int exp) {
    if (exp == 0)
        return 1;
    if (exp % 2 == 0) {
        long long half = power(base, exp / 2);
        return half * half;
    }
    return base * power(base, exp - 1);
}
  \end{lstlisting}
  \begin{itemize}
    \item Timp: $O(\log n)$ \quad Spațiu: $O(\log n)$ \textemdash{} stiva de apeluri
  \end{itemize}
  \begin{intrebare}
    De ce nu se scrie direct
    \lstinline[keepspaces=true]|return power(base, exp / 2) * power(base, exp / 2);|?
  \end{intrebare}
  \begin{itemize}
    \item<2-> Două apeluri la fiecare nivel $\Rightarrow$ numărul de apeluri se dublează la fiecare nivel
    \item<2-> În total $O(n)$ apeluri \textemdash{} câștigul se pierde; cu \lstinline|half * half|
      se face \textbf{un singur} apel pe nivel: $O(\log n)$
  \end{itemize}
\end{frame}

% ================================================================
\section{Recapitulare}

\begin{frame}{Recapitulare \textemdash{} Complexități}
  \begin{center}
  \begin{tabular}{lll}
    \toprule
    Algoritm & Timp & Spațiu \\
    \midrule
    $\gcd$ iterativ                & $O(\log \min(a,b))$ & $O(1)$ \\
    $\gcd$ extins iterativ         & $O(\log \min(a,b))$ & $O(1)$ \\
    $\gcd$ recursiv                & $O(\log \min(a,b))$ & $O(\log \min(a,b))$ \\
    $\gcd$ extins recursiv         & $O(\log \min(a,b))$ & $O(\log \min(a,b))$ \\
    Exponențiere rapidă iterativă  & $O(\log n)$         & $O(1)$ \\
    Exponențiere rapidă recursivă  & $O(\log n)$         & $O(\log n)$ \\
    \bottomrule
  \end{tabular}
  \end{center}
\end{frame}

\begin{frame}{Probleme de antrenament}
  \begin{block}{AlgoCode}
    \begin{enumerate}
      \item \href{https://code.algolymp.com/problems/1557}{\textbf{GCD and LCM}}
            \textcolor{gray}{(AlgoCode \#1557)}
      \item \href{https://code.algolymp.com/problems/4230}{\textbf{GCD and LCM Queries}}
            \textcolor{gray}{(AlgoCode \#4230)}
      \item \href{https://code.algolymp.com/problems/4270}{\textbf{GCD Power Queries}}
            \textcolor{gray}{(AlgoCode \#4270)}
    \end{enumerate}
  \end{block}
  \begin{block}{Altele}
    \begin{enumerate}
      \item \href{https://www.pbinfo.ro/probleme/58/cmmdc}{\textbf{cmmdc}}
            \textcolor{gray}{(pbinfo.ro)}
      \item \href{https://www.pbinfo.ro/probleme/59/cmmmc}{\textbf{cmmmc}}
            \textcolor{gray}{(pbinfo.ro)}
      \item \href{https://www.infoarena.ro/problema/euclid3}{\textbf{euclid3}}
            \textcolor{gray}{(infoarena.ro)}
      \item \href{https://infoarena.ro/problema/lgput}{\textbf{lgput}}
            \textcolor{gray}{(infoarena.ro)}
    \end{enumerate}
  \end{block}
\end{frame}

\begin{frame}{Resurse}
  \begin{enumerate}
    \item \href{https://www.infoarena.ro/algoritmul-lui-euclid}{\textbf{infoarena.ro}}
          \textemdash{} Algoritmul lui Euclid (RO)
    \item \href{https://cp-algorithms.com/algebra/euclid-algorithm.html}{\textbf{cp-algorithms.com}}
          \textemdash{} Euclidean algorithm (EN)
    \item \href{https://cp-algorithms.com/algebra/extended-euclid-algorithm.html}{\textbf{cp-algorithms.com}}
          \textemdash{} Extended Euclidean algorithm (EN)
    \item \href{https://cp-algorithms.com/algebra/binary-exp.html}{\textbf{cp-algorithms.com}}
          \textemdash{} Binary exponentiation (EN)
    \item \href{https://www.pbinfo.ro/articole/3873/recursivitate}{\textbf{pbinfo.ro}}
          \textemdash{} Recursivitate (RO)
    \item \href{https://cppreference.com}{\textbf{cppreference.com}}
          \textemdash{} documentație completă C++
  \end{enumerate}
\end{frame}

\end{document}
